\documentclass[a4paper,12pt,BCOR8.25mm,headsepline,final,twoside,bibtotoc,liststotoc,fleqn]{scrbook}

\usepackage[utf8]{inputenc}
\usepackage{amsfonts} %defines the \frak and \Bbb commands and set up the fonts msam (extra math symbols A), msbm (extra math symbols B, and blackboard bold), eufm (Euler Fraktur), extra sizes of cmmib (bold math italic and bold lowercase Greek), and cmbsy (bold math symbols and bold script), for use in mathematics.
\usepackage{amstext} %defines the amsmath \text command.
\usepackage{amssymb} %defines the names of all the math symbols available with the AMS fonts collection
\usepackage{amsbsy} % defines the amsmath \boldsymbol and (poor man’s bold) \pmb commands.
\usepackage{amscd} % defines some command for easing the generation of commutative diagrams.
\usepackage{amsmath} %\align, \subequation...

\hbadness=1000 % Gejammere ueber overfull/underfull boxes einstellbar (default=1000)

\usepackage{colortbl}%Tabellenfarben
\usepackage{ae,aecompl}
%\usepackage{hyphenat}
\definecolor{black}{rgb}{0,0,0}
\usepackage[plainpages=false, pdfpagelabels, colorlinks=true, linkcolor=black, menucolor=black, urlcolor=black, citecolor=black]{hyperref}


%\usepackage{subeqn} % fuer verschachtelt numerierte Gleichungen ist in amsmath enthalten
\usepackage{url}

% define the type of the thesis: Studienarbeit/Diplomarbeit
% (uncomment the appropriate type)
\newcommand{\thethesis}[0]{%
  Diplomarbeit
}

% find all usages of \field and replace the whole expression
\newcommand{\field}[1]{%
  {\itshape \{#1\}}
}





% macros to check whether running PDFLaTeX or not
\newif\ifpdf
\ifx\pdfoutput\undefined
\pdffalse % we are not running PDFLaTeX
\else
\pdfoutput=1 % we are running PDFLaTeX
\pdftrue
\pdfpkresolution 600
\pdfimageresolution 300
\pdfinfo{
   /Author (Johannes Held)
   /Title  (Die Zugrifsschicht der laufzeitadaptierten CoBRA DB)
   /Subject ()
   /Keywords (OSGi; Komponenten; Laufzeitadaption; Framework; DBMS; SOA)
}
\fi

% command to check draft option
\makeatletter
\newcommand*{\ifoptiondraft}{%
  \expandafter
  \@if@pti@ns\expandafter{\@classoptionslist}{final}%
  \@secondoftwo{%
    \expandafter
    \@if@pti@ns\expandafter{\@classoptionslist}{draft}%
    \@firstoftwo\@secondoftwo
  }%
}
\makeatother

\ifoptiondraft{%
  \usepackage[firsttwo,bottomafter]{draftcopy}
}

% include some useful packages
\usepackage{scrtime}                          % gain access to time stamps
\usepackage{scrpage2}                         % headers and footers
\usepackage{makeidx}                          % support for makeidx
\usepackage{array}                            % better table support
\usepackage{multicol}                         % spanning columns
\usepackage{multirow}                         % spanning rows
\usepackage{microtype}

\usepackage[printonlyused]{acronym}

% include right printer driver for graphicx
\ifpdf
  %\usepackage[pdftex]{graphicx}
  \usepackage{pdfpages} % lädt Paket graphics selbständig
  \pdfcompresslevel=9
\else
  \usepackage[dvips]{graphicx}
\fi

\usepackage{subfigure}
\usepackage[ngerman,english]{babel}           % switch language
\usepackage{float}                           
\usepackage[intoc]{nomencl}                          % list of abbreviations
\usepackage{algorithmic}                      % typesetting of algorithms
\usepackage[plain,chapter]{algorithm}         % typesetting of algorithms
\usepackage{stfloats}                         % used to have footnotes at bottom of the page
\usepackage[final]{listings}                  % typesetting of code listings 

\usepackage[ngerman]{varioref}
\usepackage[german]{fancyref}

\lstset{language=C++}
\lstset{basicstyle=\ttfamily\small\mdseries}
\definecolor{darkgrey}{rgb}{0.95,0.95,0.95}
\definecolor{darkgreen}{rgb}{0.3,0.6,0.3}
\definecolor{darkred}{rgb}{0.8,0.2,0.2}
\definecolor{darkblue}{rgb}{0.1,0.15,0.85}
%\lstset{backgroundcolor=\color{darkgrey}}
\lstset{stringstyle=\color{darkred}}
\lstset{numberstyle=\color{darkgreen}}
\lstset{commentstyle=\color{darkgreen}}
\lstset{keywordstyle=\color{darkblue}}
\lstset{linewidth=\textwidth, showstringspaces=false}
\lstset{captionpos=tb}
\lstset{tabsize=1}
\lstset{breaklines=true}
\lstset{frame=tlRb}
\lstset{frameround=fftt}
\lstset{breakatwhitespace=true}
\lstset{morekeywords={String,Class,Object}}
\lstset{numbers=left}
\lstset{float=htb}
\lstset{numberstyle=\ttfamily\tiny}
\lstset{numbersep=10pt}

% define command \missing
\newcommand{\missing}[1]{\,\,\textcolor{red}{(\marginpar[\hfill!$\longrightarrow$]{$\longleftarrow$!}{\bfseries 
    Missing:}\,\emph{#1})}\,\,}
\newcommand{\note}[1]{\marginpar[#1]{#1}}
\newcommand{\code}[1]{\texttt{#1}}

% environment to typeset sub-figures
\newbox\subfigbox
\makeatletter
        \newenvironment{subfloat}
                {\def\caption##1{\gdef\subcapsave{\relax##1}}%
                 \let\subcapsave\@empty
                 \setbox\subfigbox\hbox
                         \bgroup}
                  {\egroup
                 \subfigure[\subcapsave]{\box\subfigbox}}
\makeatother

% list of abbreviations
\let\abbrev\nomenclature
\renewcommand{\nomname}{Liste der Abkürzungen}
\setlength{\nomlabelwidth}{.25\hsize}
\renewcommand{\nomlabel}[1]{#1 \dotfill}
\setlength{\nomitemsep}{-\parsep}
\makeglossary
\newcommand{\markup}[1]{\textbf{#1}}

% define new style for TOC
\makeatletter
\renewcommand{\numberline}[1]{%
        \makebox[0.9cm][l]{#1}\hspace{1mm}}
\renewcommand{\l@chapter}[2]{%
        \addvspace{2ex}%
        \pagebreak[3]%
        \noindent%
        \makebox[0pt][l]{%
        \rule[-3pt]{\textwidth}{0pt}}%
        {\large\textsf{\textbf{#1}}}\hfill#2%
        \par%
        \nopagebreak%
        \addvspace{1ex}%
}
\renewcommand{\l@section}[2]{%
        \addvspace{0.5ex}%
        \noindent\hspace{1cm}%
        #1\dotfill#2%
        \par%
        \nopagebreak[2]%
}
\renewcommand{\l@subsection}[2]{%
        \addvspace{0.2ex}%
        \noindent\hspace{2cm}%
        #1\dotfill#2%
        \par%
}       
\makeatother

% define new style for index 
\makeatletter
\newcommand*{\heading}[1]{%
        \makebox[0pt][l]{%
                \rule[-3pt]{\linewidth}{0pt}}%
        \textsf{\textbf{\Large #1}}\hfill\nopagebreak\vspace{4pt}}
\renewenvironment{theindex}{%
        \setlength{\columnseprule}{0.4pt}
        \setlength{\columnsep}{2em}
        \begin{multicols}{2}[\chapter*{\indexname}]
                \parindent\z@
                \parskip\z@ \@plus .3\p@\relax
                \let\item\@idxitem}%
        {\end{multicols}\clearpage}
\makeatother

% page break
\clubpenalty = 10000
\widowpenalty = 10000

% prepare generation of index
\makeindex

% put footnotes below floats at the bottom
\fnbelowfloat

\setlength{\parindent}{0em}


\hyphenation{Ei-gen-schaf-ten}



\begin{document}


\frontmatter
\pagenumbering{alph}

\ifpdf
  \includepdf[pages={1,{}}]{intro/Deckblatt.pdf}
\fi

\begin{titlepage}
  
  \begin{center}
    
    {\Huge \bf
Entwicklung eines Frameworks für die Verteilungsoptimierung in Publish/Subscribe Systemen auf Basis eines strukturierten P2P-Overlay Netzwerks    } 
    
    \vspace*{1cm}
    Diplomarbeit im Fach Informatik
    \vspace{2cm}
    
    {\large vorgelegt von} \\
    \vspace*{0.7cm}
    {\Large \bf Johannes Held} \\
    \vspace*{0.7cm}
    {\large geb. 26.05.1983 in Nürnberg} 
    
    \vspace{2cm}
    
    angefertigt am 

    \vspace{1cm}
    
    {\bf 
      Institut für Informatik \\
      Lehrstuhl für Informatik 6\\
      Datenmanagement \\
      Friedrich-Alexander-Universität Erlangen--Nürnberg \\
      (Prof. Dr. Klaus Meyer-Wegener)
      }
    
    \vspace{1cm}
\end{center}
\begin{tabbing}
    Betreuer: \= Univ.-Prof. Dr. Richard Lenz \\
    \> Dipl.-Inf. Thomas Fischer 
\end{tabbing}
    \vspace{1cm}
    
    
\begin{tabbing}
	Trallalalalalalalalalalla \= \kill
	Beginn der Arbeit:   \> 01.06.2010 \\
  Abgabe der Arbeit:   \> 01.12.2010
\end{tabbing}
    
  
\end{titlepage}

\clearpage{\pagestyle{empty}\cleardoublepage}


\thispagestyle{empty}
\selectlanguage{ngerman}
Ich versichere, dass ich die Arbeit ohne fremde Hilfe und ohne Benutzung anderer als der angegebenen Quellen angefertigt habe und dass die Arbeit in gleicher oder ähnlicher Form noch keiner anderen Prüfungsbehörde vorgelegen hat und von dieser als Teil einer Prüfungsleistung angenommen wurde. Alle Ausführungen, die wörtlich oder sinngemäß übernommen wurden, sind als solche gekennzeichnet.
\vspace{2cm}

Der Universität Erlangen-Nürnberg, vertreten durch die Informatik 6 (Datenmanagement), wird für Zwecke der Forschung und Lehre ein einfaches, kostenloses, zeitlich und örtlich unbeschränktes Nutzungsrecht an den Arbeitsergebnissen der Diplomarbeit einschließlich etwaiger Schutzrechte und Urheberrechte eingeräumt.

\vspace{2cm}
Erlangen, den 25.09.2010

\vspace{2cm}
Johannes Held \hfill \ 

\vspace{0,5cm}

\clearpage{\pagestyle{empty}\cleardoublepage}


\selectlanguage{ngerman}

\pagestyle{useheadings}

\setlength{\parskip}{0.7em}
\pagenumbering{Roman}
\chapter*{Abstract}
\selectlanguage{english}
\missing{Bitte die deutsche Kurzfassung lesen. Übersetzung ``in progress''.}


\selectlanguage{ngerman} 
\clearpage{\pagestyle{empty}\cleardoublepage}
\chapter*{Kurzfassung}
\missing{Problemstellung muss besser eingearbeitet werden.\\\textbf{WARUM} machen wir das. Was versprechen wir uns davon?!}

Massively Multiplayer Online Games (MMOG) oder allgemeiner Massive Multiuser Virtual Environments (MMVE) sind häufig nach dem Client-Server-Prinzip aufgebaut. Die Clients, Teilnehmer der virtuellen Welt, verbinden sich mit einem oder mehreren Servern des Anbieters, auf denen die Welt verwaltet wird. Verschiedene Entwicklungen versuchen die Kommunikation über ein p2p-Netzwerk abzuwickeln und somit die virtuelle Welt durch die einzelnen Teilnehmer verwalten zu lassen. \textbf{M}assive \textbf{M}ultiuser \textbf{E}ven\textbf{t} \textbf{I}nfra\textbf{S}tructure (M$^2$etis) will diesen Ansatz durch ein optimiertes kanalbasiertes Publish/Subscribe-System erweitern. Die vorkommenden Eventtypen des MMVE bestimmen die Kanäle des Systems. Anhand der semantischen Beschreibung der Events werden die Kanäle optimiert und bieten dadurch eine auf die virtuelle Welt optimal zugeschnittene Eventverteilung. Diese Arbeit behandelt die Anbindung des Netzwerkes und die Konzeption der zur Übersetzungszeit optimierbaren Publish/Subscribe-Komponente. Hierzu werden die Grundlagen von Publish/Subscribe-Systemen und p2p-Netzwerken beschrieben. Es wurden verschiedene strukturierte p2p-Netzwerke vorgestellt und evaluiert. Pastry wurde als geeignetes Netzwerk für M$^2$etis ausgewählt. Zur Umsetzung der prototypischen Implementierung der Publish/Subscribe-Komponente sind Methoden der modernen C++-Entwicklung wie \acf{tmp} eingesetzt worden, die es erlauben, das System anhand des zur Übersetzungszeit vorhandenen Wissens zu optimieren.


\selectlanguage{ngerman}
\clearpage{\pagestyle{plain}\cleardoublepage}

\tableofcontents
\listoffigures
\clearpage{\pagestyle{plain}\cleardoublepage}
\listoftables
\clearpage{\pagestyle{plain}\cleardoublepage}
\lstlistoflistings
\clearpage{\pagestyle{plain}\cleardoublepage}
\chapter*{Abkürzungsverzeichnis}

\pagestyle{useheadings}
\markleft{Abkürzungsverzeichnis}
\markright{Abkürzungsverzeichnis}
\addcontentsline{toc}{chapter}{Abkürzungsverzeichnis}

\vspace{\topskip}


\begin{acronym}[xxxxxxxxxxxx]
	\setlength{\itemsep}{-\parsep}
	\setlength{\itemindent}{1.5em}
%A
	\acro{api}[API]{Application Programming Interface}
	\acro{aoi}[AOI]{Area of Interest}
%B

%C	
	\vspace{\parsep} 
 \acro{cast}[CAST]{group anycast/multicast}

%D	
	\vspace{\parsep}
	 \acro{dht}[DHT]{Distributed Hashtable}
	\acro{dolr}[DOLR]{Decentralized Object Location and Routing}
%E	
%	\vspace{\parsep}
%	\acro{er}[ER-Diagramm]{Entity-Relationship-Diagramm}
%F

%G	
%	\vspace{\parsep}

%H

%I	
%	\vspace{\parsep}
%	\acro{iso}[ISO]{International Organization for Standardization}
	
%J	
%	\vspace{\parsep}
	
%K
\vspace{\parsep}
\acro{kbr}[KBR]{key-based Routing}



%L

%M
	\vspace{\parsep}
	\acro{mmog}[MMOG]{Massively Multiplayer Online Game}
	\acro{m2etis}[M$^2$etis]{\textbf{M}assive \textbf{M}ultiuser \textbf{E}ven\textbf{t} \textbf{I}nfra\textbf{S}tructure}
%N

%O	
	
%P	
	\vspace{\parsep}
	\acro{p2p}[p2p]{Peer2Peer}

%Q
%	\vspace{\parsep}
	

%R	
%	\vspace{\parsep}
	\acro{rdbms}[RDBMS]{Relationales Datenbank-Management-System}
	
%S	
%	\vspace{\parsep}
	
%T	
%	\vspace{\parsep}

%U

%V
	\vspace{\parsep}
	\acro{vast}[VAST]{Voronoi-based Adaptive Scalable Transfer}
	\acro{von}[VON]{Voronoi-based Overlay Network}

%W

%X	
%	\vspace{\parsep}

%Y

%Z

\end{acronym}

\clearpage{\pagestyle{plain}\cleardoublepage}

\mainmatter

% CONTENT START

\chapter{Einleitung}
\label{chap:einleitung}

Die Spielewelt verändert sich von Einzelspieler bzw. rundenbasierten Gruppenspielen an einem PC dank der fortschreitenden Übertragungstechniken zu im Netzwerk gespielten \ac{mmog}.

Um was geht's hier:\\
Framework auf Basis eines strukurierten p2p-overlay Netwerkes zur Verteilungsoprimierung von Events in Publish/Subscribe-Systemen.

\section{Motivation}
\begin{itemize*}
\item dezidierte Server
\item Latenz
\item Ausfälle
\item Spielfluss
\end{itemize*}

\begin{itemize*}
\item vielfältige Ansätze \cite{Bharambe2008Donnybrook} %donnybrook, VAST
\item dezentral $\rightarrow$ Overlay Netzwerk
\item Pub/Sub Systeme \cite{Knutsson2004Peertopeer, Triebel2008Peertopeer} %peer2peer support
\end{itemize*}

Also im Grunde ein paar \emph{Motivationen} zusammenschreiben. :-)

\section{Aufbau der Arbeit}
Zuerst Grundlagen mit kurzer Begriffserklärung, Einführung in \ac{m2etis},  p2p-Netzwerke und Eventsystem anhand Publish/Subscribe und dessen mögliche Umsetzung in p2p-Netzwerken. In \Fref{chap:evaluation_p2p} werden verschiedene p2p-overlay Netzwerke vorgestellt und ihn ihren Arbeitsweisen einander gegenübergestellt.

\chapter{Grundlagen}
\label{chap:grundlagen}

Zur Umsetzung eines solchen Systemes werden in diesem Kapitel die Grundlagen zu \ac{p2p}-Netzwerken und Publish/Subscribe-Systemen erarbeitet. Naturgemäß werden in einem \ac{p2p}-System die einzelnen Nachrichten nicht mehr über einen vom Spielehersteller kontrollierten Server gesendet. Knoten im Netzwerk leiten Nachrichten (und damit Events) weiter und könnten daraus in betrügerischer Absicht einen Vorteil ziehen. \Fref[plain]{chap:grundlagen:cheating} geht dieser Frage nach und zeigt einen Ausschnitt der aktuellen Forschung zur Betrugsverhinderung und -aufdeckung bei \ac{mmog} in dezentralen Systemen.

Zuvor zeigt ein Szenario eines \ac{mmog} verschiedene Aktionen in einer virtuellen Welt und beschreibt auftretende Eventtypen und deren Klassifikation anhand verschiedener Dimensionen.

\section{Szenario}
\label{chap:grundlagen:szenario}
Angenommen, wir sehen einem Spieler eines typischen \ac{mmog} über die Schulter. Er wird seinen Avatar durch die Welt bewegen und mit anderen Spielern interagieren. Ein einfaches Quest zu Spielbeginn ist beispielsweise die Suche nach einem, von einem Gegner bewachten, Gegenstand. Unser Spieler muss eine Blume finden, den Gegner besiegen und die Blume aufnehmen. Diese wird zum Abschluss des Quests dem Auftraggeber gebracht und der Avatar berichtet darüber in der Gilde (ein Zusammenschluss von Spielern).\\
In diesem kleinen Beispielszenario finden sich unterschiedliche Eventtypen wie \enquote{Bewegung}, \enquote{Sprechen}, \enquote{Kampf} und \enquote{Gegenstand aufnehmen}. Die Bewegung des Avatars durch die Welt ist nur für diejenigen Spieler interessant, die diesen Avatar sehen können. Das Gespräch interessiert jedoch nur die Spieler, deren Avatare in diesem Gespräche beteiligt sind. Die einzelnen Events vom Typ \enquote{Kampf} müssen hingegen synchronisiert werden um einen fairen Schlagabtausch zu gewährleisten. Das Event \enquote{Blume aufgenommen} muss zudem abgespeichert werden um die Blume im Inventar des Avatars abzulegen und über eine Spielesession hinweg aktuell zu halten. Dieser Event ist jedoch nur für diesen Benutzer interessant.

\subsection{Semantische Beschreibung von Eventtypen}

\missing{Jetzt was sind Eventtypen, Events und die Dimensionen zur Klassifizierung von Eventtypen?}

 identifizieren und anhand verschiedner Dimensionen wie ``Kontext'', ``Persitenz'', ``Synchronität'' oder ``Validität'' klassifizieren \cite{Fischer2010Event}.




Im folgenden Kapitel werden die Grundlagen von \ac{p2p}-Netzwerken und deren grundlegende Unterteilung erläutert.

\section{p2p-Netzwerke}
\label{chap:grundlagen:p2p}

IP-Multicast ist eine Möglichkeit ein verteiltes \ac{p2p}-System aufzubauen \cite{Deering1990Multicast}. \ac{p2p}-Netzwerke setzen meist auf Overlay-Netwerke\index{Netzwerk!Overlay} auf, da IP-Multicast nicht weit verbreitet ist\footnote{Es erfordert spezielle Router und entsprechende Infrastruktur.}. Sie sind ein logischer Aufsatz auf einem bestehenden Netzwerk und abstrahieren von diesem durch einen eigenen Adressraum. Knoten im Netzwerk können benachbart sein ohne eine direkte physikalische Verbindung zu haben. Ein Overlay-Netzwerk bietet neben dem eigenen Adressraum auch Funktionalität zum Versand, Empfang und Routing von Nachrichten \cite{Tannenbaum2003}. Overlay-Netzwerke sind vielschichtiger Natur. So sind z.B. der EMail-Dienst mit seinem eigenen Namensraum, Twitter oder Facebook als Overlay-Netzwerke zu bezeichnen. p2p-Netzwerke\index{Netzwerk!p2p} selbst beschreiben den Verbund von Knoten, sogenannten Peers, die miteinander gleichberechtigt kommunizieren können \cite{Steinmetz2005}. p2p-Netzwerke sind vom unterliegenden physischem Netzwerk unabhängig, da sie auf Overlay-Netzwerke aufbauen. Obwohl Overlay-Netzwerke und \ac{p2p}-Netzwerke verschiedene Netzwerkarten sind, werden diese jedoch häufig in einem Atemzug genannt. In vielen zitierten Arbeiten wird ebenfalls von \emph{p2p overlay network} gesprochen; dies drückt die Anforderung an solch ein Netzwerk aus:\\
Ein p2p Overlay-Netzwerk ist ein vom physischen Netzwerk (räumlich) unabhängiger Verbund aus gleichberechtigten Knoten, die miteinander kommunizieren können.

Die erste Generation dieser Netzwerke wurde häufig zum Austausch von Dateien verwendet und wird daher vielfach mit dem Begriff \emph{Filesharing} assoziiert. p2p-Netzwerke werden jedoch auch als Kommunikationsnetze genutzt \cite{Darlagiannis2006Peertopeer}. Der einfache Netzaufbau, die Fehlertoleranz bei ausfallenden Knoten sowie der Wegfall teurer Server lässt diese Netzwerkart auch für Computerspiele interessant werden \cite{Knutsson2004Peertopeer, Triebel2008Peertopeer}. Roussopoulos stellt in \cite{Roussopoulos20032} einen allgemeinen Entscheidungsbaum für und wider den Einsatz von p2p-Netzwerken zur Verfügung. Auch für Sensornetzwerke ist der Aufbau eines p2p-Netzwerkes von Vorteil \cite{MuneebAliandKoenLangendoen2007Case}. Es können wertvolle Ressourcen (z.B. Batterie) besser genutzt werden \cite{Sioutas2009Building}, da die Kommunikation zwischen den Knoten weniger Sendeleistung benötigt als es eine Verbindungen zu einem -- eventuell weit entfernten -- Hauptrechner benötigen würde. Solche Netzwerke können beispielsweise zur Erkennung von Waldbränden genutzt werden. Sensorknoten werden über dem zu überwachenden Gebiet abgeworfen und sammeln Daten wie Temperatur und Luftfeuchtigkeit. Verbunden in einem großen p2p-Netzwerk werden die Daten untereinander ausgetauscht. Zur dezentralen Verwaltung und Abfrage genügt es an einem Sensorknoten die Daten abzufragen.

p2p-Netzwerke lassen sich grundsätzlich als \emph{unstrukturiert} oder \emph{strukturiert} klassifizieren \cite{Steinmetz2005, Lua2005Survey} und in Generationen einteilen \cite{Bo2003PeertoPeer}:
\begin{itemize*}
	\item \emph{1G} unstrukturierte Netzwerke
	\item \emph{2G} strukturierte Netzwerke
	\item \emph{3G} strukturierte Netzwerke mit Fokus auf Anonymität, Authentifizierung, Schutz vor Zensur und verschiedenen Rollen der einzelnen Knoten
\end{itemize*}

\subsection{Unstrukturierte Netzwerke}
Unstrukturierte Netzwerke\index{Netzwerk!Overlay!unstrukturiert} zeichnen sich dadurch aus, dass alle Informationen und Dateien durch Suchalgorithmen gefunden werden müssen \cite{Lv2002}. \\
Ein Client tritt dem Netzwerk bei und stellt seine Suchanfrage in das Netz. Werden entsprechende Peers gefunden die diese Suchanfrage beantworten können, werden zum Transfer Direktverbindungen aufgebaut. Aufgrund der einfachen (meist textbasierten) Struktur der Suchanfragen können diese einen großen Wertebereich abdecken.

Im Wesentlichen lassen sich unstrukturierte Netzwerke in die Typen \emph{zentralisiert} und \emph{dezentralisiert} einteilen.

\paragraph{zentralisiert} Knoten im Netz melden eine Liste der verfügbaren Dateien an bekannte Hauptrechner. Suchanfragen werden ebenfalls an diese Rechner gerichtet. Der Suchende erhält eine Liste von potentiellen Peers, die Dateien seiner Suchanfrage entsprechend anbieten. Diese Dateien werden über direkte Verbindungen zwischen den Peers übertragen. Die einzelnen Hauptrechner stellen bei dieser Art von System einen \emph{single point of failure} dar \cite{Eberspaecher2005}.

\paragraph{dezentralisiert} In dezentralen Netzen gibt es keine bekannten Hauptrechner. Damit sich ein neuer Knoten in das Netz einklinken kann, muss diesem mindestens ein bestehender Knoten im Netzwerk bekannt sein. Über diesen Peer tauscht der neue Knoten Informationen aus und baut eine Nachbarschaft auf. Damit Dateien gefunden werden können, wird die Suchanfrage an alle Nachbarn geschickt, sprich die Suchanfrage wird durch das Netzwerk geflutet. Diese senden sie weiter an ihre Nachbarn.\\
Die Suchanfrage kann z.B. mit einer Anzahl an Hops oder einer TTL\footnote{Time to live} in ihrer Reichweite eingegrenzt werden. Potentielle Zyklen müssen bei dieser Art von Suche aufgelöst werden \cite{Lv2002}. Sind Peers gefunden, wird ebenfalls eine direkte Verbindung zur Datenübertragung aufgebaut. 

Beispiele für solche Netzwerke sind Napster (zentralisiert), Gnutella oder BitTorrent (dezentralisiert).

\subsection{Strukturierte Netzwerke}
Strukturierte Netzwerke\index{Netzwerk!Overlay!strukturiert} der zweiten Generation sind oft dezentraler Natur. Ein Datensatz muss nicht gesucht werden, da anhand der Struktur des Netzes die zuständigen Knoten berechnet werden können. Daten können ebenfalls via Direktverbindung übertragen werden, meist wird jedoch das Netzwerk selbst zum Versand genutzt. Hierbei werden die in Nachrichten gepackte Datensätze über verschiedene Peers geroutet. Der dezentralen Art dieser Netze ist geschuldet, dass ein neuer Knoten mindestens einen Peer aus dem Netzwerk kennen muss. Viele Systeme gehen davon aus, dass der neue Knoten aus einer Liste von Peers, den ihm nächst gelegenem Knoten\footnote{Nähe im Sinne von Latenz beziehungsweise räumlicher Nähe} wählen kann und über diesen den Eintritt in das Netz anstößt.

Strukturierte Netzwerke nutzen häufig die Technik der \acf{dht}\footnote{Siehe \Fref{chap:dht} im Anhang}, da diese als Grundlage einer Vielzahl von unterschiedlichen Anwendungen dienen kann \cite{Wehrle2005, Ghodsi2006AlgorithmsDHT}. Grundsätzlich ist jedem Knoten und jedem Datensatz ein Schlüssel aus dem Schlüsselraum des Netzwerkes zugeordnet. Jeder Knoten ist dabei für einen Bereich aus dem Schlüsselraum \emph{zuständig}, dieser Bereich wird netzwerkabhängig ermittelt\footnote{Bespielsweise anhand der numerischen Nähe von Schlüsseln oder der Reihenfolge der Schlüssel.}. Für einen zu Datensatz ist damit der zuständige Knoten bekannt und das zu Grunde liegende Overlay-Netzwerk kann zur Datenübertragung genutzt werden.\\
Hier wird deutlich, dass diese Netzwerke eher der Kommunikation dienen als dem Austausch von Dateien. Darauf aufbauend gibt es unterschiedliche Systeme die sich hinsichtlich Organisation, Routing, Ein- und Austritt von Knoten und dem Verhalten im Fehlerfall unterscheiden \cite{Goetz2005, Lua2005Survey}. Beispiele für solche Netzwerke sind Chord \cite{Hosseini2007Survey}, Pastry \cite{Rowstron2001}, Tapestry \cite{Zhao2001Tapestry,Zhao2004Tapestry} oder CAN \cite{Ratnasamy2001Scalable}. 

Für \ac{m2etis} sind strukturierte Netzwerke besser geeignet, da sie einerseits als Kommunikationsmedium besser nutzbar sind als unstrukturierte p2p-Systeme bei denen meist ein massenhafter Versand der Nachricht über alle Nachbarn (\enquote{flooding}) genutzt wird. Dies widerspricht klar einer optimerten Eventverteilung im gesamten System. Weiterhin lassen sich Datensätze und Zuständigkeiten in strukturierten p2p-Netzwerken einfach berechnen, was auch die Entwicklung der Publish/Subscribe-Komponente begünstigt. Netzwerke der dritten Generation, wie zum Beispiel GNUnet \cite{Bennett2002GNet}, werden in dieser Arbeit nicht näher betrachtet. Die zusätzlich angebotene Funktionalität wird nicht benötigt und würde ob ihrer Komplexität das zu entwickelnde System aufblähen. Daher werden Chord, Pastry und CAN in \Fref[plain]{chap:evaluation_p2p}, das sich der Auswahl eines geeigneten p2p-Netzwerkes zur Verwendung mit \ac{m2etis} widmet, beschrieben und miteinander verglichen.




\subsection{Generische KBR-API}
\label{chap:grundlagen:api}
Verschiedene Typen und Implementierungen von p2p-Netzwerke haben sehr unterschiedliche Schnittstellen zur Programmierung obwohl die Anwendungsfälle meist gleicher Art sind. Dabek moniert diese unterschiedlichen Schnittstellen der verschiedenen strukturierten p2p-Netzwerke die häufig auf dem Prinzip des \ac{kbr} agieren \cite{Dabek2003Towards}. Dies mache es aus Entwicklersicht schwer, vom Netzwerk zu abstrahieren und dieses gegebenenfalls zu wechseln. Er untersucht verschiedene strukturierte p2p-Netzwerke, identifiziert ein Set von Funktionen die in zwei Teile aufgeteilt sind: \enquote{Routing messages} und \enquote{Routing state access}.\\
Erstere umfassen drei Methoden, von denen zwei \enquote{Upcalls} sind. Ein Upcall ist im Grunde eine Methode, ein Callback der Applikation die vom Netzwerk aufgerufen wird. \emph{route} ermöglicht das Senden einer Nachricht. Dabei kann ein Hinweis an das Netzwerk übergeben werden, über welchen Knoten die Nachricht als nächstes geroutet werden soll. Der Upcall \emph{forward} wird auf jedem Knoten aufgerufen, der eine Nachricht weiterleitet. Als Parameter werden der Schlüssel, die Nachricht und der nächste Routingknoten übergeben. Alle Parameter können verändert werden und der Nachrichtenversand kann auch terminiert werden. Auf dem eigentlichen Empfänger der Nachricht wird zusätzlich nach forward der Upcall \emph{deliver} aufgerufen. Als Parameter werden der Schlüssel und die Nachricht übergeben.\\
\enquote{Routing state access} umfasst vier Methoden und einen Upcall. Dieser, \emph{update}, informiert über Knoten die das Netzwerk betreten oder es verlassen. Um eine Liste möglicher Knoten die als möglicher nächster Routinghop in Frage kommen, wird die Methode \emph{local\_lookup} genutzt. \emph{neighborSet} liefert eine Liste der Nachbarn des aktuellen Knotens. Soll ein Datensatz auf mehreren Knoten gespeichert werden, ist die Methode \emph{replicaSet} nützlich. Diese liefert für einen gegebenen Schlüssel eine Liste aller Knoten, die für den Datensatz geeignet sind. Um den Bereich zu ermitteln, für den der eigene Knoten zuständig ist, wird die Methode \emph{range} angeboten.

Dabek zeigt das dieses \acf{api} ausreichend ist um darauf aufbauend verschiedene Anwendungen wie Publish/Subscribe zu implementieren und beschreibt ebenfalls wie einige Systeme -- namentlich CAN, Chord, Pastry und Tapestry -- an diese \ac{api} anpassbar sind.


\section{Betrug}\index{Betrug}
\label{chap:grundlagen:cheating}
Betrug und Betrugsverhinderung bzw. Betrugsvermeidung sind nicht Thema dieser Arbeit. Das zu entwickelnde Framework bietet genug Freiheiten betrugsresistente Publish/Subscribe-Algorithmen einzusetzen oder entsprechende Systeme auf das Framework zu setzen.\\
Dennoch soll dieses Kapitel einen kleinen Überlick über das Thema und die aktuelle Forschung dazu geben.

Betrügerisches Verhalten (Cheating) ist in p2p-Systemen einfacher als in Client/Server-Systemen, da die Kommunikation nicht zwingend über einen (vom Betreiber kontrollierten) Server läuft. Mit gefälschten Nachrichten für Statusmeldungen, wie Positionsänderungen, oder zur Entscheidungsfindung über die Reihenfolge von Aktionen können andere Knoten massiv benachteiligt werden. Das Wissen über das zugrunde liegende System wird genutzt um gezielt wichtige Aspekte des Spieles zu manipulieren.

Webb \cite{Webb2007Cheating} gibt eine Übersicht über Betrug in Netzwerkspielen und geht hier auch auf die Unterschiede zwischen Client/Server-Systemen und p2p-Systemen ein. Cheating wird in vier grundlegende Bereiche eingeteilt: \emph{Game level cheats}, \emph{Application level cheats}, \emph{Protocol level cheats} und \emph{Infrastructure level cheats}. Mögliche Verfahren wie \emph{Lockstep} \cite{Baughman2007}, das die Spielzeit in Runden einteilt, werden vorgestellt und ausgewertet.

\paragraph{Beispiel}
Als Beispiel für einen \emph{Protocol level cheat} wird das Ausnutzen des \emph{Dead-Reckoning}-Verfahrens\index{Dead-Reckoning} \cite{Pantel2002} beschrieben. Dead-Reckoning kaschiert die Latenz im Netzwerk und bietet dem Spieler durch vorberechnete Aktionen der anderen Spieler einen konstanten Spielfluss. Hierbei akzeptiert das System eine gewisse Anzahl an ausgefallenen/verlorenen Updates eines Knotens bevor unterschiedliche Aktionen aus Spielsicht durchgeführt werden.\\
Hält ein Spieler nun eigene Updates zurück, können in der Zeit bis zum geforderten Update die Informationen der anderen Spieler ausgewertet werden. Das Spiel kann somit zu eigenen Gunsten beeinflusst werden. Algorithmen die diese Problematik angehen, werden entwickelt \cite{Aggarwal2005}.

Kabus \cite{Kabus2005Addressing} beschreibt grundsätzliche Techniken die in p2p-Systemen genutzt werden um Betrug aufzudecken beziehungsweise zu verhindern. Beispielsweise können für eine Konsensus\index{Betrug!Konsensus} über Spieleraktionen zufällig Knoten gewählt werden. Diese validieren die Aktionen und können effektiv Betrug verhindern. Natürlich müssen solche Verfahren mit einem erhören Aufwand an Nachrichten teuer erkauft werden.\\
Eine weitere Möglichkeit der Betrugsverhinderung ist das Einsetzen von \emph{trusted hardware}\index{Betrug!trusted hardware} wie es bei Spielkonsolen häufig eingesetzt wird. Hier verhindert die Hardware Spielveränderungen oder den Start von externen Betrugsprogramme, die beispielsweise die Avatarsteuerung übernehmen. Dennoch besitzen auch solche Systeme Schwachstellen und können hintergangen werden.\\
Zur Betrugsaufdeckung führt Kabus die Prüfung von Logdateien an.

Auf weitere zahlreiche Arbeiten zur Betrugsverhinderung beziehungsweise -aufdeckung wird verwiesen \cite{Ferretti2008Cheating, Gauthierdickey2004Low, Kabus2007Design, Dautermann2007, Kabus2009, Castro2002Secure}.


Nach der Einführung in \ac{p2p}-Netzwerken, befasst sich das nächste Kapitel mit Publish/Subscribe-Systeme, deren Unterteilung sowie möglichen Umsetzung in \ac{p2p}-Netzwerken. 

\section{Verteilte Publish/Subscribe-Systeme}
\label{chap:grundlagen:pubsub}
Konzeptionell lassen sich Publish/Subscribe-Systeme als eventbasierte Systeme betrachten. Auf Grund ihres Aufbaus und der Skalierung  in orthogonalen Dimensionen\index{Publish/Subscribe!orthogonale Dimensionen} \enquote{Raum}, \enquote{Zeit} sowie \enquote{Synchronisation} eignen sich diese gut zur Verteilung von Events in dezentralen Umgebungen \cite{PatrickTh2003Many, Cugola2002Using}.

\begin{figure}[htbp]
\centering
\includegraphics{grafics/pubsub_black_box.pdf}
\caption{Schema eines Publish/Subscribe-Systems}
\label{fig:pubsub_black_box}
\end{figure}

Publisher und Subscriber werden durch das Event-System voneinander getrennt, wie es in \Fref{fig:pubsub_black_box} dargestellt ist.  Publisher und Subscriber sind räumlich voneinander getrennt. Ein Publisher übergibt die Nachricht an das System und hält weder direkte Verbindung mit den Subscribern noch muss der Publisher alle Subscriber kennen. Diese Trennung bezieht sich nicht nur auf verschiedene Komponenten einer Applikation, sondern kann auch über Applikations- oder gar Rechnergrenzen gehen. Die Zeitliche Trennung beschreibt, dass sich ein Subscriber am System anmelden kann obwohl kein Publisher vorhanden ist, analog können Nachrichten publiziert werden ohne dass Empfänger eingeschrieben sind. Je nach Implementierung können Nachrichten zwischengespeichert werden um diese neuen Subscribern zu zustellen. Bei einem Fernaufrufsystem wie \emph{remote proceduce call (RPC)} \cite{Birrell1984Implementing} ist dies nicht möglich, da die Gegenseite existieren muss. Das Senden einer Nachricht ist für den Publisher nicht blockierend und Subscriber warten zudem nicht aktiv auf neue Nachrichten, sondern werden meist per Callback über neue Nachrichten informiert. Damit wird die Verarbeitung vom Event-System aus nebenläufig getriggert.

Diese Arbeit beschäftigt sich ausschließlich mit dezentralen Publish/Subscribe-Sys\-temen, denn \ac{m2etis} zielt darauf ab, die Rechner der Nutzer in einem p2p-Netzwerk zu verbinden und darauf aufbauend die Events zu verteilen. Viele der Grundlagen in diesem Kapitel gelten sowohl für klassische zentrale als auch dezentrale Publish/Subscribe-Systeme, allerdings müssen im verteilten Fall die Verwaltungsinformationen ebenfalls dezentral auf allen Knoten gespeichert, beziehungsweise geeignete Verteilungsalgorithmen gefunden werden. Somit relativiert sich die Dimension der räumlichen Trennung, da Publisher wie Subscriber Teil des Eventsystems sind.

Banerjee vergleicht verschiedene Arten zum Aufbau solch eines Multicast-Systemes. \enquote{mesh-first} beschreibt den expliziten Aufbau des Netzwerkes. Die Peers untereinander verändern ihre Verbindungen aufgrund bestimmter Metriken und können auch Netzwerkpartitionen beheben sind somit selbst für das Netzwerk zuständig. Der \enquote{implizte} Ansatz beschreibt Publish/Subscribe-Systeme, die auf einem Overlaynetzwerk aufsetzen und dessen Routingalgorithmus indirekt die Verteilungsstruktur bestimmt \cite{Banerjee2001Comparative}. Ein Beispiel hierfür ist Scribe, das in \Fref{chap:related:scribe} beschrieben wird.

Fiege befasst sich näher mit dem Aspekt der Sicherheit und des Vertrauens zwischen Sender, Empfänger und dem Verteilungsystem \cite{FiegeSecurity}. Behnel stellt verschiedene Aspekte von \enquote{Quality of Service} auf verschiedenen Ebenen eines Publish/Sub\-scribe-Systems vor. Beispielsweise \enquote{Latenz}, \enquote{Bandbreite}, \enquote{Zustellgarantien} für Nachrichten auf Netzwerkebene oder \enquote{Reihenfolge}, \enquote{Validität} oder \enquote{Authentifizierung} von Nachrichten auf Verteilungsebene. Er beschreibt das Verhalten einiger Publish/Subscribe-Systeme hinsichtlich der beschriebenen Aspekte \cite{BeFiMu2006PubSubQoS}. 

Grundsätzlich lassen sich Publish/Subscribe-Systeme in zwei Varianten einteilen: \emph{kanalbasiert}\index{Publish/Subscribe!kanalbasiert} und \emph{filterbasiert}\index{Publish/Subscribe!filterbasiert} \cite{Liu2003Survey}. In kanalbasierten Systemen werden die Nachrichten einzelnen Kategorien zugeordnet. Subscriber können sich für Nachrichten dieser Kategorien anmelden und bekommen diese zugestellt. Filterbasierte Systeme haben diese Einteilung nicht, stattdessen sind Nachrichten typisiert (zum Beispiel nur einfache Datentypen und Zeichenketten) und mit einem Wertebereich versehen. Bei der Anmeldung kann ein Prädikat zur Filterung angegeben werden. Der Knoten empfängt nun nur gefilterte, auf das Prädikat passende Nachrichten.

Verbindet man die Filterung von Nachrichten mit einem kanalbasierten Ansatz gelangt man zu einem \emph{hybriden} System\index{Publish/Subscribe!hybrid}: Einer Anmeldung an einem Kanal kann ein Prädikat übergeben werden. Beispielsweise wird eine Anmeldung am Kanal für Bewegungsnachrichten über ein Gebiet eingeschränkt. Die dezentrale Filterung ist jedoch nur möglich, wenn die Nutzdaten vom System lesbar oder mit filterbaren Metainformationen angereichert sind. Zudem müssen die Prädikate im logisch aufgebauten Verteilungssystem bekanntgemacht werden, damit Nachrichten frühzeitig bei der Verteilung gefiltert werden können. Beispielsweise kann der Eventtyp \emph{Gildennachricht}\footnote{vergleiche \Fref[plain]{chap:grundlagen:szenario}} auf einem filterbaren Kanal abgebildet werden. Als Prädikat kann die Gildenzugehörigkeit des Avatars oder eine Liste der Gildenmitglieder, von denen Nachrichten erwünscht sind, angeben werden.\\
Das von \ac{m2etis} zur Verfügung gestellte kanalbasierte Publish/Subscribe-System kann pro Kanal mit einer eigenen Filterungskomponente versehen werden und somit als hybrides System genutzt werden; dies wird in \Fref{chap:konzeption_pubsub} beschrieben.

Ein prominenter Vertreter verteilter, kanalbasierter Systeme ist Scribe, dessen Funktionsweise im nächsten Abschnitt beschrieben wird.

\subsection[Umsetzung eines kanalbasieren Systemes]{Umsetzung eines kanalbasieren Systemes am Beispiel von Scribe}
\label{chap:related:scribe}
Eine Umsetzung von Publish/Subscribe-Systemen in verteilen Systemen, ist der Aufbau eines Multicast-Trees\index{Multicast-Tree}, d.h. eines durch die Knoten im Netz gebildeten Baumes in dem die Nachrichten verteilt werden. Hierbei wird pro Kanal ein eigener Multicast-Tree aufgebaut. Am Algorithmus von Scribe \cite{Castro2002Scribe}wird diese Struktur beschrieben.

Scribe basiert auf dem strukturierten Overlay-Netzwerk Pastry \cite{Rowstron2001} und erzeugt einen vom Subscriber zum Publischer aufgebauten Baum \emph{reverse path forwarding tree} \cite{Dalal1978}.

\begin{figure}[htbp]
\centering
\resizebox{\textwidth}{!}{%
\includegraphics{grafics/multicast_tree.pdf}}
\caption{Schema eines Multicast-Trees}
\label{fig:multicast_tree}
\end{figure}

\Fref{fig:multicast_tree} zeigt ein Netzwerk mit den sechs Knoten A-F. Die Verbindungen der Knoten werden durch dünne schwarze Linien dargestellt. Beispielsweise hat Knoten C Verbindungen zu A, B, D und F.\\
Der Multicast-Tree benötigt einen Knoten, der die Wurzel (im Folgenden \emph{Root-Knoten} genannt) darstellt. Aus Hashwert des Kanalnamens wird ein Schlüssel berechnet. Derjenige Knoten, der aufgrund der Netzwerkmetrik für diesen Schlüssel zuständig ist, wird Root-Knoten des Kanals. Im abgebildeten Falle ist dies Knoten A.\\
Weiterhin hält jeder Knoten eine Liste bei ihm angemeldeter Knoten. In der Abbildung wird diese Liste durch geschweifte Klammern nach der Knotenbezeichnung dargestellt.

\paragraph{Subscribe}
Knoten F sendet eine \emph{subscribe}-Nachricht an A. Diese Nachrichten sind in der Grafik durch gebogene gestrichelte schwarze Verbindungslinien mit Pfeil dargestellt. Das Netzwerk würde diese Nachricht über Knoten D und C an A routen. Knoten D lässt die Nachricht terminieren und trägt F in die Liste der Subscriber ein. Knoten D sendet nun selbst eine subscribe-Nachricht an A. C, über den die Nachricht geroutet wird, terminiert diese, trägt D in die Liste ein und sendet selbst eine subscribe-Nachricht an A. A erhält nun diese Nachricht und trägt C in die Liste ein. Damit sind nun insgesamt drei Nachrichten verschickt worden.\\
Wenn sich Knoten E für den Kanal einschreibt, wird die subscribe-Nachricht an A über den Knoten C geleitet. Dieser terminiert die Nachricht und fügt E der Liste hinzu. Da C selbst angemeldet ist, muss keine weitere Nachricht versendet werden.

Scribe fordert periodische Anmeldungen zur Erhöhung der Fehlertoleranz. Ist ein Knoten ausgefallen, routet das Netzwerk die Nachrichten über andere Knoten. Damit kann der Multicast-Tree wieder aufgebaut werden.

\paragraph{Unsubscribe}
Der Austritt aus einem Kanal erfolgt ähnlich zur Anmeldung. Die Nachricht läuft nur bis zum nächsten Knoten und terminiert dort. Der Knoten entfernt den Sender der Nachricht aus seiner Liste und sendet selbst nur eine \emph{unsubscribe}-Nachricht, wenn die Liste leer ist und er selbst nicht angemeldet ist.

\paragraph{Publish}
In \Fref{fig:multicast_tree} möchte Knoten B eine Nachricht im Kanal publizieren. B sendet darauf eine Nachricht an den Root-Knoten A, da dieser für diesen Kanal zuständig ist (gebogene türkise Linie). Nun sendet A diese Nachricht an alle Knoten in seiner Liste (gerader türkise Linie mit Pfeil). Dies ist in der Abbildung nur Knoten C. Dieser sendet sie weiter an D und E. E gibt diese Nachricht direkt an die Applikation weiter, während D die Nachricht an F schicken muss.


Hierbei ist klar ersichtlich, dass zusätzliche Nachrichten verteilt werden müssen, wenn Knoten F eine Nachricht im Kanal publizieren möchte. Diese Nachricht muss erst von Knoten F zu Knoten A wandern, damit A diese Nachricht wieder über die anderen Knoten zurücksendet. Optimierte Versionen dieses Algorithmus können hier ansetzen und zu publizierende Nachrichten nicht mehr an den Knoten senden, der ihnen diese Nachricht geschickt hat. So würde C die Nachricht nur noch an E weiterleiten.

Bayeux \cite{Zhuang2001} ist ein ähnliches System, jedoch auf Basis des Overlay-Netzwerkes Tapestry \cite{Zhao2004Tapestry}. Tapestry entspricht auch der generischen API, somit stellt dies keinen Unterschied zu Pastry dar. Im Gegensatz zu Scribe, wird bei Bayeux der Multicast-Tree vom Root-Knoten aus aufgebaut. Aufgrund der unterliegenden Routingstruktur des genutzten Overlay-Netzwerkes können sich diese Pfade unterscheiden.


Nach diesem Einblick einer möglichen Umsetzung eines kanalbasierten Publish/Sub\-scribe-Systems gibt der kommende Abschnitt am Beispiel von Mercury eine Vorstellung davon, wie filterbasierte Systeme\index{Publish/Subscribe!filterbasiert} in einem dezentralen Netzwerk implementiert sein können.

\subsection[Umsetzung eines filterbasierten Systemes]{Umsetzung eines filterbasierten Systemes am Beispiel von Mercury}
\label{chap:related:mercury}
Zur besseren Vorstellung einer Umsetzung für filterbasierte Publish/Subscribe-Systeme\index{Publish/Subscribe!filterbasiert} wird im folgenden Kapitel Mercury \cite{Bharambe2004Mercury} vorgestellt. Obwohl \ac{m2etis} ein kanalbasiertes Publish/Subscribe-System darstellt \cite{Fischer2010a}, ist es sinnvoll eine möglich Umsetzung eines filterbasierten Systems zu beschreiben um die grundlegenden Unterschiede der Systeme genauer auszuarbeiten. 

Im System gibt es eine Menge an Attributen, die ihrerseits einen definierten Wertebereich haben. Jedes Attribut wird durch einen eigenen Verbund aus Knoten, den sogenannten \emph{Hub}, bearbeitet. Der Wertebereich ist dabei nicht zwingend symmetrisch auf die Knoten verteilt.

\paragraph{Subscribe}
Eine Subscription $S$ ist ein Tupel aus Filterbedingungen über die Attribute (z.B. $S := (5 < x <= 20; y = 15)$) sowie Kontaktinformationen des Knotens. $S$ wird an einen beliebigen Knoten eines Hubs gesendet, der für das Attribut aus der Filterbedingung mit der größten Selektivität zuständig ist. Im Beispiel ist dies Attribut $y$. Im Hub wird $S$ nun zu dem Knoten weitergereicht, der den Wertebereich der Filterung abdeckt. Dort wird $S$ in einer Liste gespeichert.

\paragraph{Publish}
Eine Publikation $P$ ist ebenfalls ein Tupel mit bestimmten Werten der Attribute (z.B. $P := (x = 10; y = 0)$). $P$ wird an \emph{alle} Hubs gesendet und dort zum zuständigen Knoten weitergereicht. Dieser prüft nun die Liste der gespeicherten Subscriptions gegen die neue Publikation. Stimmen beide überein, so wird $P$ an den eingeschriebenen Knoten weitergeleitet.

Mirinae ist ebenfalls ein filterbasiertes Publish/Subscribe-System, stellt den Wertebereich eines Attributes jedoch als Hyperwürfel dar. Eine automatische Anpassung dieser Aufteilung ermöglicht eine schnelle Anpassung der Routingtabelle und damit einen kurzen Weg für die Nachrichten \cite{Choi2005Mirinae}.


\subsection{VON}
\label{chap:related:von}
\ac{von} ist in seinen Grundzügen stark unterschiedlich zu den bisher vorgestellen Umsetzung. \ac{von} nutzt das \ac{p2p}-Netzwerk nicht nur als Kommumnikationsmedium, sondern nutzt dessen Struktur auch als Verteilungsstruktur es Publish/Subscribe-Systems \cite{Hu2006VON}. VON zielt auf die Verteilungsoptimierung von Events zur Positonsänderun, muss allerdings über Applikationswissen verfügen: Die Position des Spielers. \ac{vast} \cite{Backhaus2007Voronoibased} greift das Konzept von \ac{von} auf und testet eine Implementierung auf OpenSIM \cite{Baumgart2007OverSim}.

\begin{figure}[htbp]
\centering
\includegraphics{grafics/voronoi_von_backhaus.pdf}
\caption{Struktur eines VON-Netzwerkes (aus \cite{Backhaus2007Voronoibased})}
\label{fig:von}
\end{figure}

\Fref{fig:von} zeigt einen Aufbau eines VON-Netzwerks. Die Spielwelt wird anhand der Position der einzelnen Knoten in Voronoi-Diagramme unterteilt. Jeder Knoten hat damit einen eigenen Bereich und hält Verbindungen zu seinen Nachbarn und kennt angrenzende Nachbarn im Bereich seiner \ac{aoi}. Anmeldungen im Publish/Subscribe-System sind implizit, denn Publikationen, also Positionsänderungen, werden von einem Knoten an alle direkt angrenzenden Nachbarn (``enclosing neighbor'' in \Fref{fig:von}) gesendet. Um das System konsistent zu halten werden dabei auch Informationen über ``boundary neighbors'' ausgetauscht.


Nach den Grundlagen von \ac{p2p}-Netzwerken, Publish/Subscribe-Systemen und Einblicken in verschiedene Umsetzungen, beschäftigt sich das nächste Kapitel mit der Evaluation dreier \ac{p2p}-Netzwerke m ein geeignetes System als Netzwerk für \ac{m2etis} auswählen.


\chapter{Evaluation strukturierter p2p Overlay-Netzwerke}
\label{chap:evaluation_p2p}

Dieses Kapitel bietet einen Überblick über einige p2p-Netzwerke und evaluiert diese anhand gestellter Anforderungen. Diese Evaluation beeinflusst die Entscheidung für ein Overlay-Netzwerk, das die Grundlage des zu entwickelnden verteilte Publish/Subscribe-System bildet. Diese Evaluation bedient sich zahlreicher Arbeiten, die sich alleine dem Vergleich dieser Netzwerke widmen \cite{Lua2005Survey, Goetz2005, Li2004Comparing, Darlagiannis2006Peertopeer, Castro2002Secure, Bo2003PeertoPeer} und geht auch auf ihre Nutzbarkeit als Basis für \emph{Application level multicast} sprich ein Publish/Subscribe-System ein \cite{Hosseini2007Survey, Fahmy2007, Castro2003Evaluation, Ratnasamy2001}.

Zuvor müssen jedoch die Anforderungen von \acp{mmve} an solche Systeme identifiziert werden. Diese Anforderungen werden bei der Auswahl der Netzwerkes in Betracht gezogen.

\section{Anforderungen an p2p-Netzwerke}

\paragraph{Geringe Latenz} Schnelle Reaktionszeiten und Nachrichtenübermittlung sind bei \ac{mmog} unverzichtbar. Ebenfalls müssen größere Nachrichten schnell übertragen werden, damit der Informationsfluss zur korrekten Darstellung der virtuellen Umgebung nicht behindert wird. Dies lässt sich beispielsweise anhand der Anzahl der Hops beim Nachrichtenversand messen, ist aber letzlich abhängig von der zur Verfügung stehenden Bandbreite jedes einzelnen Knotens.

\paragraph{Skalierbarkeit und Fehlertoleranz bei Knotenausfall} Selbst bei einer großen Anzahl an Knoten soll das Netz nicht kollabieren. Hierbei ist es auch wichtig, dass Knoten keine erforderliche Durchschnittszeit im Netz integriert sein müssen. Es kann davon ausgegangen sein, dass ein durchschnittlicher Spieler eines \acp{mmog} längere Zeit im Spiel angemeldet ist. Li untersucht in \cite{Li2004Comparing} wie sich p2p-Netzwerke bei großen Fluktuationen verhalten.

\paragraph{Kommunikation über das Netzwerk} Das Netzwerk soll nicht nur das schnelle Auffinden von Peers ermöglichen, sondern auch einen Transport der Nachricht (Routing) durch das Netzwerk selbst bereitstellen.

\paragraph{Eingriff in Routingentscheidungen} Applikationswissen hilft auch beim Eingriff in das Routing des Netzes. So können Knoten bevorzugt zur Weiterleitung einer Nachricht ausgewählt werden. Diese Knoten zeichnen sich beispielsweise durch eine große Bandbreite oder spezielle Applikationsmetriken\footnote{Bsp: Spieler befindet sich in der selben Stadt} aus.

\paragraph{Verfügbarkeit als C/C++-Bibliothek} Da der Prototyp dieser Arbeit in C++ entwickelt wird, muss das Netzwerk als C/C++-Bibliothek verfügbar sein. So kann das Netzwerk einfach angesprochen werden ohne dass kostenintensive Brücken zwischen beispielsweise Java und C++ geschlagen werden müssen. Da zudem betriebssystemübergreifend entwickelt und getestet wird, ist außerdem ein zur Verfügung stehender Quellcode vorteilhaft.

\paragraph{Anpassbarkeit an generische KBR-API}
Das gewählte System sollte mit der von Dabek beschriebenen generische KBR-API\footnote{siehe \Fref{chap:grundlagen:api}} kompatibel sein.

\section[Evaluation dreier p2p-Netzwerke]{Evaluation der der Netzwerke Chord, Pastry/Tapestry und CAN}
In diesem Kapitel werden die vier bekannten Systeme Chord \cite{Stoica2003}, Pastry \cite{Rowstron2001}, Tapestry \cite{Zhao2001Tapestry,Zhao2004Tapestry} und CAN \cite{Ratnasamy2001Scalable} miteinander verglichen. Die ersten drei sind in ihrem Aufbau ähnlich (der Schlüsselraum ist auf einem Ring verteilt) und unterscheiden sich in der Art des Routings. CAN hingegen bildet den Schlüsselraum auf ein d-dimensionales kartesisches Koordinatensystem ab. Alle vier Systeme sind laut \cite{Dabek2003Towards} an die generische \ac{api} anpassbar.

\subsection{Chord}
\label{chap:evaluation_chord}

\subsubsection{Aufbau und Struktur}
Chord \cite{Stoica2003} legt die $l$ bit wertigen Schlüssel (meist Zahlen im Bereich $[0,2^l-1]$) auf einem eindimensionalen Ring modulo $2^l$ im Uhrzeigersinn an. Jedem Knoten und jedem Datensatz ist ein eindeutiger Schlüssel zugewiesen, diese werden \emph{ID} und \emph{key} genannt. Ein Datensatz $X$ ist dem Knoten zugewiesen, dessen ID größer gleich dem key ist. Dieser Knoten wird Nachfolger von X, \emph{SUCC(X)}, genannt. Analog dazu gibt es auch einen Vorgänger von X, \emph{PRED(X)}.

\begin{figure}[htbp]
\centering
\includegraphics{grafics/chord_key_space.pdf}
\caption{Schlüsselraum für Chord mit sechs Knoten ($Nx$) und fünf Daten ($Kx$)}
\label{fig:chord_key_space}
\end{figure}

Damit ist ein Knoten für alle Daten zuständig, die -- bildlich gesehen -- im Ring gegen den Uhrzeigersinn vor ihm liegen. In \Fref{fig:chord_key_space} ist dies mit $l=6$ für sechs Knoten und fünf Datenpunkten gezeigt. Knoten 14 (N14) ist für den Datensatz mit Schlüssel 10 (K10) zuständig. Knoten 32 ist für K16 und K24 verantwortlich. K51 ist bei N51 zu finden. Aufgrund der Ringstruktur ist N1 für K55 zuständig. Die gestrichelten Pfeile stellen die Einträge der sogenannten Fingertabelle für Knoten $N1$ dar. 

\subsubsection{Routing}
Bei Chord besitzt jeder Knoten eine Verbindung zu seinem direkten Vorgänger und seinem direkten Nachfolger. Eine Nachricht wird von jedem Knoten so lange an seinen Nachfolger geschickt, bis sie zum zuständigen Knoten gelangt. Bei einer \texttt{lookup(x)}-Nachricht\footnote{Suche für key x Knoten N, so dass gilt: $N = SUCC(x)$.} prüft jeder involvierte Knoten A, ob sein Nachfolger für den Schlüssel zuständig ist, d.h. $ID_A < x \le SUCC(A)$. Ist dies der Fall, so sendet Knoten A die Antwort $SUCC(A)$ zurück. Bei normalem Nachrichtenaustausch wird die Nachricht zur Behandlung an den zuständigen Knoten weitergeleitet.

Da dies eine sehr ineffizientes  Routing darstellt, pflegt jeder Knoten eine Fingertabelle. Die maximal $l$ Einträge in dieser Tabelle zeigen auf andere Knoten im Ring, so dass der Eintrag in Zeile $i$ von Knoten $n$ denjenigen Knoten enthält, dessen ID mindestens um $2^{i-1}$ größer ist als $ID_n$.\\
In \Fref{fig:chord_key_space} ist die Fingertabelle von Knoten $N1$ mit ID$_{N1} = 1$ dargestellt. Die Einträge berechnen sich wie folgt: In der ersten Zeile (Idx 0) steht der Knoten, dessen ID um $2^{1-1} = 1$ größer ist, als die ID$_{N1}$. Dies ist $SUCC(2)$ und zeigt auf Knoten $N5$. In der dritten Zeile (Idx 2) findet sich der Knoten dessen ID mindestens um $2^{3-1} = 4$ größer ist als ID$_{N1}$, also $SUCC(5)$ und damit ebenfalls auf Knoten $N5$. Der vierte Eintrag verweist demnach auf $SUCC(9) = N14$. Analog dazu berechnen sich die restlichen Einträge.

Über diese Fingertabelle können Nachrichten eine größere Strecke auf dem Ring überbrücken und die Routingzeit wird stark verkürzt. Da die IDs in der Tabelle exponentiell zur Basis $2$ ansteigen, halbiert sich die Distanz zum Ziel. Damit besitzt das Routing eine Komplexität von $O(log N)$.

\subsubsection{Nachbarschaft}
Die Nachbarschaft ist bei Chord begrenzt. Jeder Knoten hat eine Verbindung zu seinem Vorgänger sowie Nachfolger auf dem Ring und hält Einträge in der Fingertabelle vor. In die Routingentscheidungen kann nicht direkt eingegriffen werden.

\subsubsection{Eintritt und Austritt (Fehlerfall) von Knoten}
Bei Chord kann die ID für einen neuen Knoten $n$ frei aus dem Wertebereich des Schlüsselraums gewählt werden, es muss lediglich ein Knoten $b$ im System bekannt sein. $n$ routet eine \texttt{lookup(n)}-Nachricht via $b$ und erfährt somit, wer sein Nachfolger auf dem Ring ist. In gleicher Weise vervollständigt $n$ seine Fingertabelle. Weiterhin teilt $n$ seinem Nachfolger mit, dass $n$ nun sein neuer Vorgänger ist.

Zur Stabilisierung und Vervollständigung der Routinginformationen arbeitet jeder Knoten im Ring periodisch die Funktion \texttt{stabilize} ab. Jeder Knoten $a$ fragt seinen Nachfolger $n$ nach dessen Vorgänger $s$. Wenn $s$ ungleich $a$ ist, muss Knoten $s$ neu in den Ring eingetreten sein. $a$ informiert den neuen Knoten $s$, dass er sein Vorgänger ist und ändert selbst den eigenen Nachfolger auf $s$ ab. $s$ kennt nun den Bereich seiner Zuständigkeit $[ID_a, ID_s[$, fordert diese Daten von seinem Nachfolger an und weist diesen auf die Zuständigkeitsänderung hin.\\
Die Aktualisierung der Fingertabelle \texttt{fix\_fingers} wird ebenfalls auf jedem Knoten periodisch angestoßen. Für einen zufällig gewählten Eintrag wird überprüft, ob dieser noch aktuell ist.

Die verzögerte Aktualisierung hat keinen großen Einfluss auf die Korrektheit oder Geschwindigkeit des Routings, da eine fehlerhafte Nachrichtenzustellung über die Nachfolger- beziehungsweise Vorgängerverbindungen des empfangenden Knotens weitergeleitet wird.

\begin{figure}[htbp]
\centering
\resizebox{\textwidth}{!}{%
\includegraphics{grafics/chord_new_node.pdf}}
\caption{Schlüsselraum von Chord nach Ankunft von Knoten $N20$}
\label{fig:chord_new_node}
\end{figure}

\Fref[plain]{fig:chord_new_node} verdeutlicht den Neueintritt von Knoten $N20$. Die Änderungen sind in Grün dargestellt: $N1$ passt seine Fingertabelle an und $N20$ ist für $K16$ zuständig.

Der Ausfall von Knoten wird über \emph{Timeouts} ermittelt. Im Falle eines Timeouts wird die Nachricht an den besten bekannten Vorgänger des ausgefallenen Knotens weitergeleitet. Im schlimmsten Falle ist dies der Nachfolger des sendenden Knotens. Daraus wird ersichtlich, dass ein valider Nachfolger notwendig ist. Aus diesem Grund hält jeder Knoten eine Liste von möglichen Nachfolgern vor, die während \texttt{stabilize} erstellt werden. Für fehlerhafte Knoten in der Fingertabelle kann \texttt{fix\_fingers} explizit aufgerufen werden.
Knotenausfall bedeutet nicht nur einen Ausfall des Knotens, sondern bedingt, dass die dort gespeicherten Daten nicht mehr erreichbar sind.\\
Verlässt ein Knoten das Netz, so beeinflusst dies das System nicht. Jedoch ist es effizienter, wenn ein verlassender Knoten seinem Vorgänger und Nachfolger dies mitteilt, die Verbindungen angepasst werden und die Datensätze explizit übertragen werden.

\subsection{Pastry/Tapestry}
\label{chap:evaluation_pastry}

\subsubsection{Aufbau und Struktur}
Pastry \cite{Rowstron2001} und Tapestry \cite{Zhao2001Tapestry,Zhao2004Tapestry} sind sich sehr ähnlich, da beide auf Plaxtons Arbeit \cite{Plaxton1997Accessing} aufbauen. Daher wird auf Tapestry nicht näher eingegangen, da die Unterschiede für das Konzept dieser Systeme nicht relevant sind.

Pastry besitzt wie Chord einen $l$ bit-wertigen Schlüsselraum, dabei werden Schlüssel als Zahlen zur Basis $2^b$ dargestellt, wobei die Wahl von $b$ einen Einfluss auf das Routing hat, da dieser die Größe der Routingtabelle beeinflusst. Ein Datensatz ist dem Knoten zugewiesen, dessen ID den kleinsten Abstand zum Schlüsselwert des Datensatzes hat.\\
\Fref{fig:pastry_key_space} zeigt dies beispielhaft für die gleichen sechs Knoten und fünf Datensätze wie in \Fref[plain]{fig:chord_key_space}. Im Unterschied zu Chord ist jedoch Knoten $N14$ für $K16$ und Knoten $N54$ für $K55$ zuständig.


\begin{figure}[htbp]
\centering
\includegraphics{grafics/pastry_key_space.pdf}
\caption{Schlüsselraum für Pastry mit sechs Knoten ($Nx$) und fünf Daten ($Kx$)}
\label{fig:pastry_key_space}
\end{figure}

\subsubsection{Routing}
Jeder Knoten verwaltet neben den ihm zugeteilten Datensätzen drei Strukturen, die dem Routing dienen. Diese sind das \emph{leaf set} mit Einträgen zu Knoten, die im Schlüsselraum benachbart sind, das \emph{neighborhood set} mit Einträgen zu Knoten, die aus Netzwerksicht nahe liegen, und die Routingtabelle selbst. In \Fref{fig:pastry_routing_table} sind nur die Schlüssel, nicht aber Kontaktinformation, wie IP-Adresse und Port dargestellt.

\begin{figure}[htbp]
\centering
\includegraphics{grafics/pastry_routing_table.pdf}
\caption{Routing table, leaf set und neighborhood set bei Pastry}
\label{fig:pastry_routing_table}
\end{figure}

Die Routingtabelle besteht aus $\frac{l}{b}$ Reihen mit je $2^b -1$ Einträgen. \Fref{fig:pastry_routing_table} zeigt dies beispielhaft für Knoten $103220$ mit $l=12, b=2$ (nach \cite{Goetz2005}). Einträge in Zeile $i$ haben einen Präfix der Länge $i$ mit dem Knoten $103220$ gemein. Die Übereinstimmungen sind in der Abbildung fett hervorgehoben. Ist kein passender Knoten bekannt, wird das entsprechende Feld nicht ausgefüllt. Damit hat die Routingtabelle Ähnlichkeiten zur Fingertabelle bei Chord. Ein Knoten hat ungenaues Wissen über entfernte Knoten. Der Detailgrad an Routinginformationen erhöht sich pro Zeile in der Routingtabelle. Sind im System wenig Knoten vorhanden, sind die letzten Reihen der Routingtabelle spärlich belegt. Im Durchschnitt sind bei $n$ Knoten im System nur $log_{2^b} n$ Einträge ausgefüllt.\\
Bei der Belegung der Routingtabelle werden bei gleichem Präfix diejenigen Knoten gewählt, die aus Netzwerksicht näher sind.

Das leaf set enthält die $l$ numerisch nahen Knoten, $\frac{l}{2}$ davon sind kleiner und $\frac{l}{2}$ größer als der aktuelle Knoten. Neben Informationen zu Routingentscheidungen wird es zur Reparatur genutzt, sollten nahe gelegene Knoten ausfallen.

Das eigentliche Routing unterscheidet zwei Fälle: Zuerst prüft der Knoten, ob der Schlüssel $k$ im Bereich seines leaf sets ist. Ist dies der Fall, wird die Nachricht an den entsprechenden Knoten gesendet. Ist dieser Knoten für den Schlüssel zuständig, endet das Routing. Fällt $k$ nicht in den Bereich des leaf sets, wird die Nachricht via Routingtabelle an einen weiter entfernten Knoten gesendet. Hierzu wird ein Eintrag gewählt, der eine größere beziehungsweise die größte Präfixübereinstimmung mit $k$ hat. Existiert kein solcher Eintrag, wird die Nachricht an den numerisch nächsten Knoten (zu $k$) mit gleicher Präfixübereinstimmung geschickt.\\
Da Nachrichten immer an Knoten mit einer größeren Übereinstimmung oder an nähere Knoten mit gleicher Präfixübereinstimmung gesendet werden, können keine Zyklen auftreten.

Dadurch verringert sich die Anzahl der Knoten mit längeren Präfixübereinstimmungen in jedem Schritt um mindestens den Faktor $2^b$. Somit hat das Routing eine Komplexität von $O(log_{2^b} N)$.

\subsubsection{Nachbarschaft}
Das neighborhood set (siehe \Fref[plain]{fig:pastry_routing_table}) enthält $|m|$ Knoten, die aus Netzwerksicht nahe sind. Obwohl es im Routing keine Rolle spielt, kann es dazu genutzt werden, um geeignete Knoten zu finden. Da die Größe des leaf sets ebenfalls wählbar ist, können hier auch vermehrt nahe Knoten platziert werden.

\subsubsection{Eintritt und Austritt (Fehlerfall) von Knoten}
Einem neuen Knoten $n$ wird von Applikationsseite ein frei wählbarer Schlüssel gegeben. Meist berechnet sich dieser Schlüssel anhand des Hashwerts über die IP oder seines öffentlichen Namens. Weiterhin geht das System davon aus, dass $n$ aus einer Liste bekannter Knoten denjenigen Knoten $x$ wählen kann, der aus Netzwerksicht nahe ist. Von diesem Knoten kann das neighborhood set kopiert werden. Zum Aufbau der Routingtabelle und des leaf sets lässt $n$ via $x$ eine \texttt{JOIN}-Nachricht an einen numerisch nahen Schlüssel zu $n$ routen. Diese Nachricht gelangt schließlich zu Knoten $c$, von dem das leaf set (die Liste der $l$ numerisch nahen Knoten) kopiert werden kann. Alle Knoten, die diese \texttt{JOIN}-Nachricht weiterleiten, senden ihre Routingtabelle an $n$. Für jeden Routinghop kopiert sich $n$ die entsprechende Zeile aus der Routingtabelle, da ausgehend von keiner Präfixübereinstimmung mit dem nahen Knoten $x$, jeder weitere Hop eine größere Präfixübereinstimmung bringt.\\
Im Gegensatz zur nachträglichen Aktualisierung bei Chord wird hier die gesamte Routinginformation an alle bekannten Knoten gesendet. Der neue Knoten ist nun im Netzwerk bekannt und erreichbar.

Ausgefallene Knoten werden ebenfalls anhand von Timeouts beim Routing entdeckt. Da die Einträge des neighborhood sets nicht im Routing involviert sind, müssen diese periodisch geprüft werden. Fehlerhafte Einträge in der Routingtabelle können über einen anderen Eintrag mit gleicher Präfixübereinstimmung kompensiert werden, müssen aber entfernt werden, um ein stabiles und sicheres Routing zu gewährleisten. Hierzu können von benachbarten Einträgen Routinginformationen angefordert werden, um die entstandene Lücke zu füllen. Ein fehlerhafter Eintrag im leaf set oder neighborhood set kann auf ähnliche Weise repariert werden: Hier werden Informationen von anderen, im leaf set oder neighborhood set eingetragenenen, Knoten angefordert.

Der Austritt eines Knotens wird vom System wie ein Fehlerfall behandelt. Um die Datenintegrität zu gewährleisten und um unnötigen Nachrichtenversand im System zu vermeiden, sollte die Applikation den Austritt eines Knotens jedoch speziell behandeln.


\subsection{CAN}
\label{chap:evaluation_can}

\subsubsection{Aufbau und Struktur}
Der Schlüsselraum bei CAN \cite{Ratnasamy2001Scalable} ist ein d-dimensionaler Torus. Die Schlüssel werden als d-Tupel (zum Beispiel $(x,y)$ für $d=2$) dargestellt. Wie bei Pastry ist der numerisch nächst gelegene Knoten für einen Datensatz zuständig. Der Schlüsselraum ist in nicht überlappende Zonen eingeteilt und hat eine feste Größe. Jeder Knoten \emph{besitzt} eine solche Zone, über deren Ausmaß er definiert ist. Er ist für alle Daten zuständig, deren Schlüssel in dieser Zone liegt. Ein Schlüssel wird, wie bei \ac{dht} üblich, über eine segmentierte Hashfunktion berechnet\footnote{Siehe \Fref{chap:dht} im Anhang}. Jedes Segment bildet dabei eine Dimension ab.

\Fref{fig:can_key_space} zeigt einen zweidimensionalen Schlüsselraum mit den drei Knoten A, B und C und fünf Datesätzenn, die als schwarze Punkte dargestellt sind. Der Schlüsselraum ist komplett auf die drei Knoten aufgeteilt, wobei A für den Bereich $(0, .5)-(1, 1)$, B für $(0, 0)-(.5, .5)$ und C für $(.5, 0)-(1, .5)$ zuständig ist.

\begin{figure}[htbp]
\centering
\includegraphics{grafics/can_key_space.pdf}
\caption{Zweidimensionaler Schlüsselraum für CAN}
\label{fig:can_key_space}
\end{figure}

\subsubsection{Routing} 
Die gespeicherte Routinginformation ist bei CAN am geringsten: Jeder Knoten speichert lediglich seine Nachbarn, das sind Knoten deren Zonen angrenzend sind, ab. Über die Zoneninformation jedes Nachbarn wird nun das Routing bestimmt. Eine Nachricht wird über den Knoten geschickt, der aufgrund seiner Zoneninformation dem Ziel am nächsten ist.

In \Fref{fig:can_routing} hat Knoten N1 die vier Knoten N2, N3, N4 und N5 als Nachbarn. Knoten N4 hingegen ist nur mit N1 und N5 benachbart. Eine Nachricht von $N1$ zu $K$ wird beispielsweise via Knoten $N5$ geroutet. Eine alternative Route via $N2$ ist gestrichelt dargestellt.

\begin{figure}[htbp]
\centering
\includegraphics{grafics/can_routing.pdf}
\caption{Routing und Nachbarschaft bei CAN}
\label{fig:can_routing}
\end{figure}

Die durchschnittliche Länge eines Pfades bei $d$ Dimensionen und $n$ Knoten ist $O(d\cdot n^\frac{1}{d})$. Sind die Zonen gleich aufgeteilt, so besitzt jeder Knoten max. $2d$ Nachbarn und die Anzahl an Hops verringert sich auf $O(\frac{4}{d}\cdot n^\frac{1}{d})$.

\subsubsection{Nachbarschaft}
Die Nachbarschaft von CAN wurde bereits im vorigen Abschnitt behandelt. Dank deren Einfachheit ist der Ein- und Austritt von Knoten zwar besonders einfach und tangiert nur wenige Knoten im Netz, jedoch kann auf die Nähe aus Netzwerksicht nur beim Eintritt eines Knotens bei der Wahl seiner Koordinaten Rücksicht genommen werden.

Die Erhöhung der Dimension bedingt eine größere Nachbarschaft und damit, neben einem kürzeren Routing, auch die Möglichkeit, mehr Knoten in der Nachbarschaft anzusiedeln. Zusätzlich können in CAN sogenannte \emph{Realitäten} genutzt werden. Dies sind verschiedene CAN-Netzwerke mit unterschiedlichen Hashfunktionen zur Berechnung der Schlüssel. Jeder Knoten und jeder Datensatz ist in jedem dieser Netzwerke vertreten und besitzt einen unterschiedlichen Schlüssel pro Netzwerk. In einem System mit $r$ Realitäten muss ein Knoten demnach $r$ verschiedene Zonen und Nachbarschaften verwalten. Eine Nachricht wird bei jedem Routinghop über die Realität verschickt, welche die kürzeste Route verspricht.

\subsubsection{Eintritt und Austritt (Fehlerfall) von Knoten}
Für den Eintritt eines neuen Knotens $n$ muss wieder ein Knoten $b$ aus dem Netz bekannt sein. Eine Koordinate im Schlüsselraum wird n zugewiesen und eine spezielle \texttt{JOIN}-Nachricht via $d$ an die gewählte Koordinate gesendet. Ist diese Nachricht über das normale Routing bei dem für diese Koordinate zuständigen Knoten $d$ angekommen, halbiert dieser seine Zone und weist eine Hälfte dem neuen Knoten $n$ zu. Die Aufteilung der Zonen erfolgt dabei anhand einer Reihenfolge der Dimensionen. Dies vereinfacht die Aufteilungs- sowie die Zusammenführungsprozedur. Letztlich kopiert $d$ alle Daten aus dieser neuen Zone zu $n$. Knoten $d$ schickt $n$ seine Nachbarschaftsinformationen und trägt diesen selbst als neuen Nachbarn ein. Knoten n teilt seine Anwesenheit sofort seinen neuen Nachbarn mit. Der Neueintritt eines Knotens ist damit auf wenige Nachrichten zwischen den Nachbarn begrenzt und beeinträchtigt das übrige Netzwerk nicht. Über periodische \texttt{UPDATE}-Nachrichten halten sich Nachbarn stets aktuell und senden ihre eigenen Nachbarschaftsinformationen an die benachbarten Knoten.

\begin{figure}[htbp]
\centering
\includegraphics{grafics/can_new_node.pdf}
\caption{Eintritt und Fehlerfall bei CAN}
\label{fig:can_new_node}
\end{figure}

\Fref{fig:can_new_node} zeigt den Eintritt für Knoten N7. Nach Teilung der Zone von N4 verändert sich die Nachbarschaft von N1: Knoten N7 kommt neu hinzu. Die Nachbarschaft von $N1$ ist dann $\{N2, N3, N4, N5, N7\}$. Auf der rechten Seite ist der Ausfall des Knotens $N2$ (rot) dargestellt. Knoten $N3$ ist nun für die vereinigte Zone zuständig.

Ausstehende \texttt{UPDATE}-Nachrichten oder Timeouts weisen auf ausgefallene Knoten hin. Werden diese zum Routen einer Nachricht benutzt, stellt dies keine direkte Beeinträchtigung des Netzes dar, da die Nachricht über einen anderen Nachbarn verschickt werden kann. So könnte N1 auch N2 zum Nachrichtenversand nutzen, wenn N5 ausgefallen wäre (vergleiche \Fref[plain]{fig:can_routing}).

Ein Knoten, der einen ausfallenden Peer entdeckt hat, startet einen Timer, dessen Dauer proportional zur Zonengröße ist. Nach Ablauf sendet der Knoten eine spezielle \texttt{TAKEOVER}-Nachricht an alle Nachbarn des ausgefallenen Knotens\footnote{Dieses Wissen ist durch vorherige \texttt{UPDATE}-Nachrichten bekannt.}. Erreicht eine \texttt{TAKEOVER}-Nachricht einen Knoten, der den Ausfall ebenfalls bemerkt hat, stoppt dieser seinen eigenen Timer, falls seine Zone größer als die des Senders ist. Andernfalls antwortet er selbst mit seiner \texttt{TAKEOVER}-Nachricht. Auf diese Weise wird der Nachbar mit der kleinsten Zone gefunden. Dieser ist nun der neue Besitzer der Zone und fügt seine beiden Zonen zusammen, wie es in \Fref{fig:can_new_node} am Beispiel von N2 und N3 ersichtlich ist. N2 (rot dargestellt) ist ausgefallen und N3 vergrößert seine Zone. N2 ist nun nicht mehr in der Nachbarschaft von N1 enthalten. Es ist möglich, dass ein Knoten Besitzer zweier Zonen wird, welche sich nicht als eine Zone darstellen lassen. Ein Beispiel hierfür ist Knoten N4, der beim Ausfall von N5 dessen Zone übernimmt. Ein Hintergrundprozess defragmentiert solche Zonen und weist Knoten gegebenenfalls neue Koordinaten zu. Eine periodische Auffrischung der Datensätze soll auch hier einen Verlust im Fehlerfall schnell kompensieren.

Möchte ein Knoten $l$ das System verlassen, so sucht er einen Nachbarn mit der kleinsten Zone und sendet diesem seine Daten. Dieser Nachbar informiert nun die alten Nachbarn von $l$ über die geänderte Nachbarschaft.


\section{Auswahl des p2p-Netzwerkes}
Nach der Vorstellung der verschiedenen Netzwerke werden diese im Folgenden verglichen um ein geeignetes System als Netzwerk für \ac{m2etis} zu finden. Jedes System hat eigene Schwächen und Stärken. So bezahlt CAN einen günstigen Ein- und Austritt der Knoten dank der kleinen Nachbarschaft mit mehr Routing Hops als beispielsweise Chord. Dieses hat durch seine Fingertabelle wiederum mehr Einfluss auf das Routing, aktualisiert diese Tabellen jedoch erst nachträglich im Hintergrund während dies bei Pastry aktiv geschieht.

\Fref{tab:evaluation_fazit} (Auszug aus \cite{Goetz2005}) listet die durchschnittlich anfallenden Kosten (Routing Hops, Größe der Routinginformation, Nachrichtenanzahl beim Ein- und Austritt) für die drei getesteten System auf. $n$ ist die Anzahl der Knoten, $b$ die Anzahl der Bits der darstellenden Basis bei Pastry und $d$ die Anzahl der Dimensionen bei CAN.

\begin{table}[htbp]
\centering
\begin{tabular}{lcccc}
\toprule
 & Routing Hops & Routinginformation & Eintritt & Austritt\\ 
 \cmidrule{2-5}
Chord & $O(\frac{1}{2}log_2~n)$ & $O(2log_2~n) $ & $ O(log_2^2 n) $ & $ O(log_2^2 n) $ \\
Pastry & $O(\frac{1}{b}log_2~n)$ & $O(\frac{1}{b} (2^b-1) log_2~n) $ & $ O(log_{2^b}~n) $ & $ O(mlog_b~n) $ \\
CAN & $O(\frac{d}{2}n^\frac{1}{d})	$ & $O(2 d) $ & $ O(\frac{d}{2}n^\frac{1}{d}) $ & $ O(2 d) $ \\
\bottomrule
\end{tabular}
\caption[Vergleich der Systeme Chord, Pastry und CAN]{Vergleich der Systeme Chord, Pastry und CAN anhand einiger Gesichtspunkte}
\label{tab:evaluation_fazit}
\end{table}

Für uniformes Routing (d.h. das Routing ist bei allen Peers gleich) sind $O(log~n)$ beziehungsweise $O(n\frac{1}{d})$ Hops für Routingtabellen der Größe $O(log~n)$ beziehungsweise $O(d)$ die asymptotischen Grenze \cite{Xu2004Fundamental}. Die vorgestellten \ac{dht}-basierten Netzwerke nutzen ein uniformes Routing.


\paragraph{Geringe Latenz und Kommunikation über das Netzwerk}
Bei Pastry/Tapestry sind die kürzesten Routen und damit auch oftmals geringste Latenz im Nachrichtenversand zu erwarten, da die Routingtabelle im Vergleich zu Chord und CAN mehr Einträge enthält und diese einfacher mit - aus Netzwerksicht - nahen Peers belegt werden kann. Bei CAN hingegen wird das langsamste Routing erwartet, da Nachrichten nur zwischen benachbarten Knoten ausgetauscht werden. Sprünge (via Fingertabellen) sind nicht vorgesehen. Die Erwartungen decken sich mit den Werten in \Fref{tab:evaluation_fazit}.

Li stellt für verschiedene Parametereinstellungen der Netzwerke deren Bandbreite und Latenz gegenüber und untersucht dabei das Verhalten bei gehäuften Ein- und Austritten von Knoten. Hier ist Chord leicht im Vorteil, da lediglich der Verweis auf den Nachfolgeknoten für das korrekte Routing erforderlich ist. Bei allen Netzwerken pendelt sich die Latenz im \emph{worst case} auf 250ms ein \cite{Li2004Comparing}.

Die Vorteile der Kommunikation bei Pastry/Tapestry überwiegen für unseren Anwendungsfall.

\paragraph{Skalierbarkeit}
CAN steht hier an erster Stelle, da Ein- und Austritte von Peers nur wenige Knoten im Netz betreffen. Auch bei vielen Ein- und Austritten leidet das Netz unter keiner großen Nachrichtenlast. Allerdings leidet die Kommunikation in großen Netzen (siehe obigen Punkt). Chord liegt auf dem letzten Platz, da das Netzwerk erst durch später aufgerufene Methoden vollkommen funktionsfähig wird (siehe \emph{Eintritt und Austritt (Fehlerfall) von Knoten}).  Pastry und Tapestry liegen auf Zweitem Platz; der Aufbau der Routingtabelle erfolgt in vielen kleinen Schritten - dafür ist ein Knoten danach ein vollwertiger Peer im System.

\paragraph{Fehlertoleranz bei Knotenausfall}
Alle Systeme finden ausgefallene Knoten durch periodisch versandte Nachrichten an benachbarte Peers oder durch Timeouts verschickter Nachrichten. Jedes System kompensiert solch einen Fall auf eigene Art und Weise. Bei CAN wird ein Knoten seine mit der verlassenen Zone verbinden, Knoten in Chord aktualisieren ihre Fingertabellen und Peers in Pastry und Tapestry versuchen durch Befragungen anderer Knoten die entstandene Lücke in der Routingtabelle zu füllen. Ein Knotenausfall hat bei CAN allerdings die geringsten Auswirkungen auf das restliche System, da nur die Nachbarknoten involviert sind. Dies ist der kleinen Routingtabelle geschuldet.

Alle Systeme fordern jedoch die Applikation zu einer periodischen Auffrischung der gespeicherten Daten auf und bieten selbst kaum Redundanz (Ausnahme Tapestry) an.


\paragraph{Bestimmung der Nachbarschaft}
Allein die Größe der Routingtabelle bedingt, dass bei Pastry und Tapestry mehr Einfluss auf die Zusammenstellung genommen werden kann. Bei CAN gibt es faktisch nur eine Entscheidung bei Eintritt in das Netz, Chord bietet über die Fingertabelle minimalen Einfluss, während bei Pastry explizit das Neighborhood Set eingesetzt wird, um eventuelle Lücken in der Routingtabelle geschickt zu besetzen.

\paragraph{Eingriff in Routingentscheidungen}
Das Routing bei CAN kann nur bedingt beeinflusst werden. Chord hingegen bietet mit seiner Fingertabelle mehr Variationsmöglichkeiten - allerdings nur für die zu überbrückende Distanz im Schlüsselraum. Pastry und Tapestry verbinden diese Variationen mit vielfältigen Nachbarschaftsoptionen.

Zusammenfassend zeigen die obigen Punkte, dass Pastry/Tapestry als Netzwerk für \ac{m2etis} gut geeignet ist. Pastry selbst ist als Javabibliothek\footnote{\url{http://www.freepastry.org}} verfügbar und die Entwicklung von Tapestry (ebenfalls in Java implementiert) wurde mit Version 2.01 eingestellt. Chimera \cite{Allen2006Chimera} ist der Nachfolger von Tapestry und vereint das Beste von Pastry und Tapestry in sich\footnote{siehe \url{http://current.cs.ucsb.edu/projects/chimera/index.html}}: 
\selectlanguage{english}
\begin{quote}
Chimera is a light-weight C implementation of a \enquote{next-generation} structured overlay that provides similar functionality as prefix-routing protocols Tapestry and Pastry.  Chimera gains simplicity and robustness from its use of Pastry's leafsets, and efficient routing from Tapestry's locality algorithms.  In addition to these properties, Chimera also provides efficient detection of node and network failures, and reroutes messages around them to maintain connectivity and throughput.  
\end{quote}
\selectlanguage{ngerman} 

Der frei verfügbare Code (veröffentlicht unter GPL), die Anpassbarkeit und die Unterstützung der Zielplattformen Linux und Windows sprechen für Chimera. Weiterhin entspricht Chimera der generischen \ac{api} und kann bei gravierenden Problemen durch ein anderes System ausgetauscht werden, ohne das restliche System zu beeinflussen. Die Konzeption des Netzwerkes ist ausschlaggebender als eigentliche Implementierung und daher wird Chimera als Netzwerk für \ac{m2etis} gewählt.

In diesem Kapitel wurden die drei Netzwerke Chord, Pastry/Tapestry und CAN verglichen und ihre unterschiedlichen Arbeitsweisen erklärt. Es zeigt sich, dass das Konzept von Pastry/Trapestry für \ac{m2etis} gut geeignet ist. Damit ist der erste Teil dieser Arbeit, die Auswahl eines geeigneten Netzwerkes, abgeschlossen. Das nächste Kapitel befasst sich nun mit der Konzeption der Publish/Subscribe-Komponente, die auf diesem Netzwerk aufsetzt.

\chapter{Konzeption des Frameworks}
\label{chap:konzeption_pubsub}
In diesem Kapitel wird die Konzeption des Frameworks zur Verteilungsoptimierung in seinen Einzelheiten erläutert.\\
Ausgehend von den in \cite{Fischer2010a} identifizierten Typen von Events und deren orthogonalen Dimensionen zur Optimierung wird die, im nächsten Abschnitt beschriebene, Problemstellung greifbar. Das Framework muss die vielfältige Anpassung einzelner Kanäle ermöglichen ohne dies mit Einbußen zur Laufzeit zu erkaufen. Ein weiterer wichtiger Fokus dieser Arbeit ist es \ac{m2etis} als ein einfach zu benutzendes kanalbasiertes Publish/Subscribe-System zu präsentieren:\\
Es muss ohne Wissen über die verschiedenen Optimierungen verwendbar sein.

\section{Umsetzung der Dimensionen}
Für die logische Umsetzung der in \cite{Fischer2010Event} identifizierten semantischen Dimensionen eines Eventtyps wird der Begriff \emph{Policy} eingeführt.\\
Eine Policy definiert die Schnittstelle für verschiedene konkrete Implementierungen (genannt Strategie) und deren Auswirkung auf die Nachrichtenverarbeitung im Publish/Subscribe-System. Die folgenden sieben Policies decken die Dimensionen ab.

\begin{description}
\item[Verteilung] bestimmt die Verteilungsart der einzelnen Events und den Aufbau des logischen Multicast-Trees, mittels dem die Nachrichten versandt werden \cite{KostasKatrinis2005}.
\item[Filterung] erlaubt es Anmeldungen Prädikate mitzugeben. Diese Policy stellt sicher, dass diese Prädikate nach oben im Multicast-Tree zusammengeführt werden und Nachrichten frühzeitig gefiltert werden können. Dies bedeutet, dass Nachrichten jeweils beim Versand durch den logischen Kopf des Multicast-Trees gefiltert werden.
\item[Zustellung] bestimmt das Kommunikationspradigma des Nachrichtenversands und leitet beispielsweise den Versand von Bestätigungen über eingegangene Nachrichten an den sendenden Knoten ein.
\item[Reihenfolge] definiert das Synchronisationskonzept eines Kanals.
\item[Persistenz] bietet die Möglichkeit der Speicherung eines Event beim Empfänger.
\item[Sicherheit] gibt eine Schnittstelle zur Nachrichtenverschlüsselung vor.
\item[Validität] prüft die ankommenden Nachrichten auf ihre Validität. Frühzeitig verworfene Nachrichten vermindern das Nachrichtenaufkommen im System stark.
\end{description}

Im Optimierungsschritt erzeugt \ac{m2etis} für jeden semantisch Typen einen optimierten Kanal der über das Publish/Subscribe-System ansprechbar ist\footnote{siehe \Fref{chap:grundlagen:aufbau_metis}}. Dieses bietet die bekannten Methoden \emph{subscribe, unsubscribe} und \emph{publish} an. Jeder Kanal ist entlang den Dimensionen derart optimiert, dass Nachrichten bestmöglich verarbeitet und über das Netzwerk verteilt werden können.\\
Das Netzwerk selbst wird über die bereits erwähnte \ac{kbr}-API angesprochen. Drei Methoden, sogenannte ``upcalls''\footnote{Technisch gesehen eine Callbackmethode der Applikation, welche vom Netzwerk aufgerufen wird.},  informieren über Ein- und Austritte von Knoten sowie über zu routende und auch ankommende Nachrichten an einem Knoten. Andere Methoden ermöglichen die Abfrage von Netzwerkinformationen \cite{Dabek2003Towards}.

Die Verbindung der Publish/Subscribe Methoden mit der KBR-API ist trivial: Jede verschickte Nachricht wird durch die Methoden \emph{forward} und \emph{deliver} des Netzwerkes verarbeitet. Strategien hingegen können weitere Methoden zum Informationsgewinn nutzen.\\
Die Verteilung der Policies auf die verschiedenen Nachrichtentypen ist hingegen interessanter und in \Fref{tab:verbindungsmatrix} dargestellt. Hierbei ist zu beachten, dass \emph{Publish}-Nachrichten den eigentlichen Events entsprechen, deren Verteilung optimiert wird, und dass \emph{Subscribe}- und \emph{Unsubscribe}-Nachrichten nur zum Aufbau und der Verwaltung des Multicast-Trees dienen. Hierzu sind nur die Policies Verteilung, Filterung und Sicherheit nötig, während für Publish-Nachrichten alle Policies involviert sind.

\begin{table}[!h]
\resizebox{\textwidth}{!}{%
\begin{tabular}{llccccccc}
\toprule
\multirow{2}{2cm}{Nachrichten\-typ} & \multirow{2}{2cm}{KBR-API Methode}	& \multicolumn{7}{c}{Policy pro Kanal} \\
\cmidrule{3-9}
			&	& Verteilung & Filterung & Zustellung & Reihenfolge & Persistenz & Sicherheit & Validität \\
\toprule 
publish	    & deliver & + & + & + & + & + & + & + \\
\midrule
\multirow{2}{*}{subscribe}	& deliver & + & + &   &   &   & + & \\
\cmidrule{2-9}
			& forward & + & + &   &   &   & + & \\
\midrule
\multirow{2}{*}{unsubscribe} & deliver & + & + &   &   &   & + & \\
\cmidrule{2-9}
			& forward & + & + &   &   &   & + & \\
\bottomrule
\end{tabular}}
\caption{Verbindungsmatrix}
\label{tab:verbindungsmatrix}
\end{table}

Die eingesetzte Verteilungsstrategie entscheidet darüber, ob Subscribe- und Unsubscribe-Nachrichten in forward und/oder deliver behandelt werden. Nur in foward ist die Nachricht durch Verteilung und Filter veränderbar und kann auch durch den Verteilungsalgorithmus terminiert werden. Durch die Behandlung aller Publish-Nachrichten -- also Events -- ausschließlich in deliver, wird sichergestellt, dass diese bei allen Empfängern ankommen und unterwegs nicht verändert oder unterbrochen werden.

Eine Ausweitung der anderen Policies auf die anderen Nachrichtentypen ist nicht nötig. Anmeldungen und Abmeldungen bedingen keiner speziellen Reihenfolge und müssen zudem auch nicht gespeichert werden. Eine spezielle Zustellungsgarantie für diese Nachrichten ist kontraproduktiv:  Anmeldungen müssen -- systembedingt -- periodisch wiederholt werden. Eine Zustellbestätigung würde die Anzahl der Verwaltungsnachrichten stark erhöhen.

Nachdem die Verbindung des Netzwerkes mit der Publish/Subscribe-API geschildert ist, wird die Reihenfolge der Policies -- bestimmend für eine effiziente Bearbeitung jeder einzelnen Nachricht -- im nächsten Kapitel über das Verarbeitungsmodell beschrieben.

\section{Verarbeitungsmodell}
Das von der Verteilungspolicy erzeugte System ist ein logischer Multicast-Tree. Es gibt einen oder mehere Knoten, im folgenden Rootknoten genannt, welche die Wurzel des Multicast-Trees darstellen. Bei Scribe\footnote{siehe \Fref{chap:related:scribe}} gibt es einen Rootknoten pro Kanal. Nachbarschaftszentrierte Algorithmen, die zum Beispiel an VON \cite{Hu2006VON} angelehnt sind, haben mehere Knoten als Wurzel. 

Beim Erstellen einer Nachricht, werden Verwaltungsinformationen der einzelnen Policies abgefragt und mit der Nachricht verschickt. Anhand dieser Informationen kann die Nachricht auf einem anderen Knoten entsprechend behandelt werden. Subscribe- und Unsubscribe-Nachrichten bestehen nur aus Verwaltungsinformationen, da sie zum Aufbau und der Verwaltung des Multicast-Trees dienen und keine Events transportieren.

Der Versand einer Nachricht ist in \Fref{fig:processing_send} dargestellt. Nachdem die Nachricht erstellt ist, wird die Verteilungspolicy nach einer Liste von Empfängern befragt. Diese sind vom Nachrichtentyp abhängig.\\
Für Subscribe- und Unsubscribe-Nachrichten sind dies oft die Rootknoten, denn diese koordinieren den Aufbau. Publish-Nachrichten werden an die Rootknoten geschickt, falls die Nachricht mit \emph{to root} gekennzeichnet ist. Dies ist der Fall, wenn ein Knoten einen Event in das System bringen möchte. Ein Event wird meist an die Rootknoten geschickt, denn die Verteilung wird in der Methode deliver ausgeführt. Hat die Nachricht die Rootknoten erreicht, wird sie mit \emph{from root} gekennzeichnet und die Verteilungsspolicy gibt als Empfänger die eingeschriebenen Knoten aus. Für jeden Eintrag dieser Liste wird nun von der Filterpolicy geprüft, ob die Nachricht auf das mit diesem Knoten verbundene Prädikat passt und an diesen verschickt werden soll.\\
Schließlich wird die Nachricht durch die anhand der Sicherheitspolicy vorgegebene Verschlüsselung kodiert und über das Netzwerk an alle Empfänger gesandt.

\begin{figure}[htbp]
\centering
\resizebox{\textwidth}{!}{%
\includegraphics{grafics/processing_send.pdf}}
\caption{Versand von Nachrichten}
\label{fig:processing_send}
\end{figure}

Werden Nachrichten in forward behandelt, müssen diese erst dekodiert werden, wie es in \Fref{fig:processing_forward} aufgezeigt ist. Die Verteilungs- und Filterpolicies können anhand der Verwaltungsinformationen ihren Zustand anpassen und wenn nötig die Nachricht ändern oder gar stoppen.

\begin{figure}[htbp]
\centering
\resizebox{\textwidth}{!}{%
\includegraphics{grafics/processing_forward.pdf}}
\caption{Verarbeitung von Nachrichten in forward}
\label{fig:processing_forward}
\end{figure}

Die Abarbeitung der Nachrichten in deliver ist komplexer als die beiden oben genannten Fälle, da hier alle Policies zusammenarbeiten wie es \Fref{fig:processing_deliver} aufzeigt. Nach der Entschlüsselung werden Subscribe- und Unsubscribe-Nachrichten ähnlich wie bei der Behandlung in forward verarbeitet. Die Policies können ihren Zustand aktualisieren.\\
Publish-Nachrichten müssen eine Validitätsprüfung bestehen, bevor entschieden wird, ob sie eine Nachricht \emph{to root}, also an die Rootknoten sind, oder nicht. Jeder Rootknoten und alle anderen Knoten auf dem Verteilungsweg, die selbst Verteilungsaufgaben übernehmen müssen (abhängig von der gewählten Verteilungspolicy), leiten nun das Verteilen der Nachricht ein. Dazu wird in einem ersten Schritt eine neue Publishnachricht \emph{from root} erstellt, alle Verwaltungsinformationen abgefragt und schließlich die Nachricht mit dem obig beschriebenen Verfahren gesendet. Nun wird an diesen Knoten geprüft ob sie selbst an diesem Kanal angemeldet sind und, falls dies zutrifft, ob sie an der Nachricht interessiert sind. Wenn nicht, endet die Bearbeitung der Nachricht. Ansonsten trifft sich der Ablaufpfad an dieser Stelle mit dem Ablaufpfad einfacher Knoten, die lediglich Subscriber sind. Die Synchronisierungspolicy ermöglicht eine Wohlordnung der Nachricht und kann diese sowohl komplett zurückhalten als auch mehrere Nachrichten zurückgeben, falls durch die aktuelle Nachricht weitere Nachrichten ``freigeschaltet'' werden. Alle Nachrichten werden nun nochmals auf ihre Validität geprüft, da zurückgehaltene Nachrichten inzwischen veraltet sein können. Für jede valide Nachricht wird eine Signalisierung der Zustellung laut Policy, zum Beispiel eine Bestätigungsnachricht zurück an den Sender, ermöglicht. Bevor die Nachrichten schließlich an die Applikation übergeben werden, können sie persistiert werden.

\begin{figure}[htbp]
\centering
\resizebox{\textwidth}{!}{%
\includegraphics{grafics/processing_deliver.pdf}}
\caption{Verarbeitung von Nachrichten in deliver}
\label{fig:processing_deliver}
\end{figure}

Der nächste Abschnitt prüft das eingeführte Verarbeitungsmodell anhand einiger Strategien auf seine Tauglichkeit. Dazu werden hauptsächlich Verteilungsstrategien genutzt, da diese die Hauptarbeit im System tragen.

\subsection*{Beispielhafte Strategien}
Für jede einzelne Policy bietet \ac{m2etis} mehrere konkrete Implementierungen an und lässt sich mit benutzerdefinierten Strategien weiter auf die Anforderungen der Applikation anpassen\footnote{siehe \Fref{fig:metis_aufbau}}. In diesem Kapitel wird das Verarbeitungsmodell anhand zweier Verteilungsalgorithmen getestet. Scribe \cite{Castro2002Scribe} erzeugt wie beschrieben einen Multicast-Tree, während der andere Algorithmus an VON \cite{Hu2006VON} angelehnt ist und seinen Fokus auf die ``in-game''-Nachbarschaft legt. Dies ist vor Allem auf Events mit einer hohen Lokalität sinnvoll. Der Einfluss solcher Events, wie zum Beispiel Positionsänderungen, wirkt sich nur in einem begrenzten Gebiet aus und muss daher nur benachbarten Knoten -- im Spiel -- mitgeteilt werden.

Zuerst wird das Verarbeitungsmodell zum Nachrichtenversand und danach der Empfang und die Verteilung von Nachrichten betrachtet.

Bei Scribe richten sich im Gegensatz zu VON die Zielknoten nach der Nachrichtenart und dem Typ des Knotens. Subscribe- und Unsubscribe-Nachrichten, wie auch Publish-Nachrichten normaler Knoten, werden immer an den Rootknoten des Channels gesendet. Hierzu kann der Hashwert des Kanalnamens berechnet und auf einen Schlüssel im Netz abgebildet werden. Verteilt der Rootknoten oder weitere Knoten auf dem Verteilungsweg eine Publish-Nachricht so gibt Scribe die Lister der an diesme Knoten eingeschriebenen Subscriber zurück.\\
Ein Algorithmus wie VON kann auf Applikationswissen zugreifen und liefert als Empfänger einer Subscribe-Nachricht die Nachbarn im Spiel sind zurück. Diese Nachbarschaftmetrik muss sich nicht alleine auf die Position im Spiel beziehen, sondern kann auch von der Mitgliedschaft in Gruppen bestimmt werden; sie ist aber in jedem Falle spezifisch für den gewählten Eventtyp. Unsubscribe-Nachrichten richten sich prinzipell an die gleichen Nachbarn -- jedoch müssen diese mit den Nachbarn zur Anmeldung abgeglichen werden. Publish-Nachrichten zur werden zuerst immer an den Knoten selbst geschickt, damit dieser in der Abarbeitung in deliver den Event an alle eingetragenen Empfänger senden kann.

Scribe verarbeitet Subscribe- und Unsubscribe-Nachrichten in forward. Jeder Knoten auf dem Routingpfad vom Sender zum Rootknoten fügt den Sender der Nachricht seiner Liste der Empfänger hinzu und ändert die Anmeldenachricht: Er trägt sich als Absender ein. Der Multicast-Tree wird also Knoten für Knoten beim Routing der Nachrichten aufgebaut. Sollte der Knoten schon angemeldet sein, kann er die Nachricht terminieren und damit effektiv die Anzahl an Nachrichten begrenzen.\\
Die von Scribe geforderte periodische Auffrischung der Anmeldung erfolgt für normal angemeldete Knoten außerhalb der Verteilungsstrategie: Die Anmeldung wird nach einem, von der Strategie bestimmten Zeitintervall, wiederholt. Ist ein Knoten aufgrund seiner Lage im Multicast-Tree angemeldet, kann er diese Anmeldung auffrischen, wenn periodische Anmeldungen zu bearbeiten sind. Hierbei ist es wichtig, dass die Filterstrategie bei allen Auffrischungen involviert ist, damit die Prädikate weiterhin im Baum zusammengeführt werden. Durch die periodischen Anmeldungen werden auch ausgefallene Knoten ausgetauscht, denn die Nachrichten werden vom Netzwerk über andere Knoten geroutet und somit wird der Multicast-Tree wieder aufgebaut.\\
Der VON-ähnliche Algorithmus bearbeitet alle Nachrichten in deliver. Der Sender der Subscribe-Nachricht wird in die Liste der Empfänger eingetragen. Eine Publish-Nachricht wird an all diese weitergeleitet. Da ein Knoten nicht bei sich selbst angemeldet ist, kann die Bearbeitung der Nachricht beendet werden.

Es ist allerdings auch möglich eine Verteilungsstratege wie VON gänzlich anders zu implementieren. Alle Anmeldungen für sind lediglich implizit und werden nicht als Nachrichten auf das Netzwerk gelegt. Statt eine Publish-Nachricht an den eigenen Knoten zu senden, wird diese allen Nachbarn im Spiel zugestellt. Dies bedeutet, dass Publish-Nachrichten vom Sender mit \emph{from root} statt \emph{to root} markiert werden.\\
Mit diesem implizitem System entfällt auch die Abmeldung am Kanal.

Diese drei vorgestellen Implementierungsansätze verschiedender Verteilungsstrategien zeigt wie mächtig das Verarbeitungsmodell ist und welchen Spielraum es verschiedenen Strategien bietet um Events anhand der angebotenen Dimensionen bestmöglich zu bearbeiten.

\chapter{Prototypische Implementierung}
\label{chap:impl}
Dieses Kapitel widmet sich der prototypischen Implementierung der im vorigen Kapitel vorgestellten Konzeption der Publish/Subscribe-Komponente von \ac{m2etis}. Bevor der Aufbau dieser Komponente beschrieben wird, wird die Entwicklungsumgebung (Programmiersprache und gewähltes \ac{p2p}-Netzwerk) beschrieben. Danach zeigt der Zugriff auf das Publish/Subscribe-System die einfache Schnittstelle zu \ac{m2etis} aus Benutzersicht. Die vom Optimierungsschritt zu erstellenden beziehungsweise anzupassenden Dateien werden im darauffolgenden Kapitel beschrieben. Details aus der Implementierung schließen dieses Kapitel ab.

\section{Wahl der Entwicklungsumgebung}
Als Programmiersprache wird C++ gewählt, da diese Sprache die gesamte Bandbreite der Softwareentwicklung anbietet und sowohl den Zugriff auf die Hardware des Computers als auch die Entwicklung von großen Systemen begünstigt. C++ legt ebenso einen Schwerpunkt auf die Sprachmittel zur Entwicklung von Bibliotheken und favorisiert dadurch verallgemeinerte Mechanismen gegenüber in die Sprache integrierten Einzellösungen für typische Problemstellungen. Eine der Stärken von C++ -- als Multiparadigmensprache -- ist dabei die Kombination von effizienter, maschinennaher Programmierung mit mächtigen Sprachmitteln zur Anwendungsentwicklung. Dabei wird in dieser Arbeit vor allem die, in \Fref[plain]{chap:impl_tmp} im Anhang beschriebene, \acf{tmp} benutzt.

Mit der Festlegung auf C++ muss auch die Implementierung des p2p-Netzwerkes nativ durch C++ ansprechbar sein. Das bedeutet, das Netzwerk muss als C- oder C++-Bibliothek vorliegen. Chimera \cite{Allen2006Chimera} ist der Nachfolger von Tapestry und wird auf der Homepage\footnote{siehe \url{http://current.cs.ucsb.edu/projects/chimera/index.html} (26.\,11.\,2010)} folgendermaßen beschrieben: 
\selectlanguage{english}
\begin{quote}
Chimera is a light-weight C implementation of a \enquote{next-generation} structured overlay that provides similar functionality as prefix-routing protocols Tapestry and Pastry.  Chimera gains simplicity and robustness from its use of Pastry's leafsets, and efficient routing from Tapestry's locality algorithms.  In addition to these properties, Chimera also provides efficient detection of node and network failures, and reroutes messages around them to maintain connectivity and throughput.  
\end{quote}
\selectlanguage{ngerman} 

Pastry\footnote{\url{http://www.freepastry.org} (26.\,11.\,2010)} und Tapestry sind als Java-Bibliotheken verfügbar. Die Entwicklung von Tapestry wurde mit Version 2.01 eingestellt. Chimeras frei verfügbarer Code (veröffentlicht unter GPL) und die Kompabilität mit der generischen KBR-API sprechen weiterhin für Chimera. Bei gravierenden Problemen kann es somit durch ein anderes Netzwerk ausgetauscht werden, ohne dass diese Änderung großen Einfluss auf die restlichen Komponenten von \ac{m2etis} hat. 


\section{Aufbau der Publish/Subscribe-Komponente}

Das UML-Klassendiagramms in \Fref[plain]{fig:uml} spiegelt den groben Aufbau der Komponenten wieder und zeigt die Verbindung der wichtigsten Klassen. Auf die Darstellung von Hilfskonstrukten wird der Übersichtlichkeit halber verzichtet. Als erster Schritt wird die Anbindung an das \ac{p2p}-Netzwerk beschrieben. Danach wird -- ausgehend von der Applikation -- die Implementierung des Publish/Subscribe-Systems erläutert. Angewandte Techniken wie \ac{tmp} oder \emph{policy based}-Design werden im Anhang dieser Arbeit beschrieben. Sie werden benötigt um den zusätzlichen Verwaltungsaufwand zur Laufzeit durch geschickte Anwendung des zur Übersetzungszeit vorhandenen Wissens zu minimieren. Das nutzbare Wissen umfasst beispielsweise die Typen der Events oder die zur Optimierung ausgewählten Strategien und deren Besonderheiten.

\texttt{P2PInterface} ist eine abstrakte Basisklasse und repräsentiert die in \Fref[plain]{chap:grundlagen:api} beschriebene \ac{kbr}-API. Aktuell wird diese durch die Klasse \texttt{ChimeraWrapper} implementiert, welche über die Klasse \texttt{ChimeraWrapperImpl} Zugriff auf das Netzwerk Chimera bietet. Hier wurde das \enquote{PIMPL}-Pattern angewandt, mit dem die eigentliche Implementierung versteckt werden kann \cite{Alexandrescu2001Modern}. Damit werden die Implementierungsdetails, wie interne Datentypen oder Methoden, vor dem Nutzer versteckt und müssen demnach nicht im zu inkludierenden Header angegeben werden. Diese können nun nach Belieben geändert werden, ohne dass der Nutzer neukompilieren muss, da sich die Schnittstelle nicht verändert. Als Nebeneffekt sinkt die Übersetzungzeit, da insgesamt weniger Header inkludiert und geparst werden müssen.

Zur Vereinfachung werden sämtliche Identifikationsdaten, zum Beispiel die Adresse eines Knotens, des Netzwerkes als \texttt{std::string} repräsentiert. Dies entbindet den Entwickler der Implementierung weiterer Wrapperklassen von komplexen Hilfskonstrukten. Eventuell muss daher in der Implementierungsklasse (in unserem Fall \texttt{ChimeraWrapperImpl}) eine zusätzliche Indirektion eingeführt werden, um den \texttt{std::""string} wieder umzusetzen. Zur eigentlichen Nachrichtenübertragung wird der Datentyp \texttt{std::vector$<$char$>$}, die objektorientierte Version eines \texttt{char*}, bereitgestellt. Neben der einfacheren Zugriffsweise bietet dies zusätzlichen Schutz beispielsweise gegen einen Pufferüberlauf.\\
Weitere Netzwerke können auf einfache Art und Weise durch entsprechende Wrapperklassen eingebunden werden.

\begin{figure}[htbp]
\centering
\resizebox{\textwidth}{!}{%
\includegraphics{grafics/uml.pdf}}
\caption{Vereinfachtes Klassendiagramm des Frameworks}
\label{fig:uml}
\end{figure}

Die \texttt{Applikation} greift auf das Publish/Subscribe-System über den Singleton \texttt{PubSub\-Sytem} zu. Beiden ist der optimierte \texttt{Channel} nur über die abstrakten Basisklassen \texttt{IChannel} beziehungsweise \texttt{IChannelOwn} zugänglich. Diese sind in \Fref[plain]{lst:interface_channel} dargestellt und bieten lediglich die drei üblichen Methoden\footnote{\texttt{subscribe}, \texttt{unsubscribe} und \texttt{publish}} für Publish/Subscribe-Systeme an. Sie verdecken dadurch die Komplexität der Klasse \texttt{Channel} vor dem Benutzer. Bei der Anmeldung an einem Kanal muss ein Funktionspointer vom Typ \texttt{app\_deliver\_func} übergeben werden. Diesem übergibt das Publish/Subscribe-System eine empfangene Nachricht und liefert damit die Events an die Applikation aus. Der Anmeldung kann ebenfalls ein Prädikat zur Filterung übergeben werden, falls die gewählten Strategien dies unterstützt. \texttt{Channel} selbst ist eine mit policy-based Design erstellte Templateklasse\footnote{siehe \Fref{chap:impl_tmp} im Anhang zur Erklärung} und bekommt die verschiedenen Strategien -- in \Fref{fig:uml} ebenfalls nicht dargestellt -- als Policies übergeben. \ac{tmp} ermöglicht es zudem, jede \texttt{Channel}-Instanz auf die gewählten Policies abzustimmen. Weiterhin wird für jeden \texttt{Channel} ein optimierter Header erzeugt, dessen Größe an Nutzlast für Verwaltungsdaten von den gewählten Policies abhängig ist. Durch diese Maßnahmen wird der Overhead zur Laufzeit stark reduziert. 

\lstinputlisting[caption={Schnittstellen der Klasse \texttt{Channel}}, label=lst:interface_channel, float=!t]{listings/interface_channel.h}

Die Klasse \texttt{PubSubController} implementiert das Interface \texttt{P2PCallback} und kann sich somit für die Callbacks des Netzwerkes registrieren. Der \texttt{PubSubController} ist die Schnittstelle des Publish/Subscribe-Systems mit dem \ac{p2p}-Netzwerk und bietet die benötigte Funktionalität für die Klasse \texttt{Channel}. Zudem werden die ankommenden Nachrichten in \texttt{deliver} über eine Queue und einen Dispatch-Thread an die einzelnen Kanäle verteilt. 

\texttt{PubSubSystemImpl} ist eine der generierten Klassen (siehe nächstes Kapitel) und kennt die verschiedenen Netzwerkwrapper (in \Fref{fig:uml} ist dies \texttt{ChimeraWrapper}). Aktuell besitzt diese für jeden optimierten \texttt{Channel} einen Pointer auf \texttt{IChannelOwn} als Klassenvariable. Auf diese hat die Klasse \texttt{PubSubSystem} als \texttt{friend class}\footnote{\texttt{friend} ist eine besondere Beziehung zwischen zwei Klassen} freien Zugriff und kann die Methodenaufrufe an der Publish/Subscribe-API an die einzelnen Instanzen weiterleiten. Für jedes genutzte Netzwerk instantiiert \texttt{PubSubSystemImpl} einen eigenen \texttt{PubSubController} und übergibt diesen an die entsprechenden Kanäle. Mit verschiedenen Instanzen der Klasse \texttt{PubSubController} können somit auch verschiedene Netzwerke angesprochen und entsprechend der Optimierungsmetriken für verschiedene Kanäle genutzt werden.

Die UML-Beschreibung der als \emph{Singleton} implementierten Klasse \texttt{PubSubSystem} -- die zur Interaktion mit \ac{m2etis} dem Nutzer zur Verfügung steht -- ist in \Fref{fig:uml_pubsubsystem} dargestellt. Zur Vereinfachung sind lediglich die \texttt{public} verfügbaren Methoden angegeben. Die Schnittstelle der Klasse erlaubt \emph{method chaining}, den einfachen Aufruf mehrerer Methoden hintereinander und erfordert meist die Angabe des Kanals -- kodiert in ein \texttt{enum} des Typs \texttt{ChannelList}, welches im nächsten Abschnitt genauer beschrieben wird.  Durch die überladene Methode \texttt{init}, wird der Eintritt des Knotens in das \ac{p2p}-Netzwerk angestoßen. Der erste Knoten baut das Netzwerk auf, während die restlichen Knoten nur über bekannte Knoten in das Netzwerk eintreten können. Nur durch Methoden der Klasse \texttt{PubSubSystem} kann sich ein Knoten für einen Kanal an- beziehungsweise abmelden. Über ein \emph{Handle} auf einen Kanal kann lediglich publiziert werden. Das Beispiel in \Fref[plain]{lst:pubsub_usage} zeigt, wie einfach diese Schnittstelle benutzt werden kann.

\begin{figure}[htbp]
\centering
\includegraphics{grafics/uml_pubsubsystem.pdf}
\caption{UML-Beschreibung der Klasse \texttt{PubSubSystem}}
\label{fig:uml_pubsubsystem}
\end{figure}

Dank der strikten Aufteilung wird die Nutzerfreundlichkeit des Frameworks aus Entwicklersicht ermöglicht, wie es das Codebeispiel im nächsten Abschnitt darlegt. Hier ist deutlich zu sehen, dass der Nutzer keinerlei Wissen über das genutzte Netzwerk noch über die zur Optimierung herangezogenen Strategien haben muss, um das System benutzen zu können.

\section{Zugriff auf M$^2$etis aus Benutzersicht}
Das \Fref[plain]{lst:pubsub_usage} zeigt den Zugriff der Applikation auf das Publish/Subscribe-System. Der Code erzeugt drei Knoten, die mit dem System arbeiten. Der erste Knoten (auf Port 7000) sendet Nachrichten auf dem Kanal \texttt{ANY\_ALL}\footnote{Der Name \texttt{ANY\_ALL} suggeriert, dass hier jegliche Art von Nachrichten an alle Knoten verteilt werden.}, für den sich die beiden anderen Knoten (Ports 7001 und 7002) anmelden. In den Zeilen 10 bis 12 ist die Empfangsfunktion definiert, welche dem System bei Anmeldungen übergeben werden muss. Diese Funktion wird in den Zeilen 18 und 21 mittels \texttt{boost::bind} an die benötigte Signatur angepasst und mit Metadaten (hier den Knotennamen) angereichert. Der erste Knoten wird in den Zeilen 15 und 16 initialisiert und ein Handle auf den Kanal \texttt{ANY\_ALL} erlangt. Die Anmeldung der Knoten erfolgt in den Zeilen 19 und 22, mit Angabe eines im Netzwerk bekannten Knoten zum Einstieg in das Netzwerk und der Empfangsfunktion. In den Zeilen 24\,-\,27 wird in einer Endlosschleife eine Nachricht über das Handle in das System gebracht und fünf Sekunden geschlafen.


\lstinputlisting[caption={Zugriff auf M$^2$etis aus Benutzersicht}, label=lst:pubsub_usage, float=!t]{listings/pubsub_usage.cpp}

Dieses Codebeispiel nutzt das angebotene \emph{method chaining} und zeigt deutlich die Kapselung des komplexen, optimierten Publish/\-Subscribe-Systems und dessen einfache und unkomplizierte Handhabbarkeit. Über \texttt{boost::bind} können beliebige Methoden an die Signatur der Empfangsfunktion angepasst werden. Daher muss die Applikation einerseits nicht von einer abstrakten Basisklasse abgeleitet sein und andererseits bleibt die eigentliche Empfangsfunktion anpassbar, wie es im Codebeispiel dargestellt ist.

Nach diesem Beispiel zur einfacher Schnittstellennutzung, wird im folgenden auf die im Optimierungsschritt\footnote{Schritt 2 in \Fref{fig:metis_aufbau}} zu erstellenden Klassen eingegangen.

\section{Anforderungen an den Optimierungsschritt}
Die Dateien und Klassen müssen in der richtigen Reihenfolge generiert und angepasst werden um die verschiedenen Kanäle zu erzeugen und mit dem Netzwerk zu verbinden. \Fref{fig:reihenfolge} zeigt diese Reihenfolge der nun vorgestellten Dateien.

\begin{figure}[!h]
\centering
\includegraphics{grafics/reihenfolge.pdf}
\caption{Erstellungsreihenfolge der Dateien im Optimierungsschritt}
\label{fig:reihenfolge}
\end{figure}


\subsubsection*{\texttt{ChannelList.h}}
Der \ac{m2etis}-Optimierer muss die Datei \texttt{ChannelList.h} komplett generieren und alle vorhandenen Kanäle -- die je einem Eventtyp entsprechen -- benennen und im \texttt{enum} eintragen, wie es \Fref{lst:channellist} darstellt. Die Einträge dienen der Auswahl des Kanals bei der Kommunikation mit dem \texttt{PubSubSytem}, wie es \Fref[plain]{lst:pubsub_usage} zeigt. Durch das \texttt{enum} wird ein zur Laufzeit einfach zu überprüfendes Mapping auf die eigentliche Implementierung des Kanal ermöglicht und verhindert einen dem System unbekannten Kanal anzufordern.

\lstinputlisting[caption={Beispiel der generierten Datei \texttt{ChannelList.h}}, label=lst:channellist, float=!t]{listings/channellist.h}


\subsubsection*{\texttt{PubSubSystemImpl.h}}
\texttt{PubSubSystemImpl} ist -- wie schon beschrieben -- die Implementierung der offentlichen Schnittstelle für Nutzer von \ac{m2etis}. Sie muss komplett generiert werden, da sie Zugriff auf alle zur Optimierung genutzten Strategien hat und damit die daraus erzeugten optimierten Instanzen der Templateklasse \texttt{Channel} als Klassenvariablen besitzt.

\Fref[plain]{lst:pubsubsystemimpl} zeigt einen Ausschnitt dieser generierten Klasse. In den Zeilen 3 und 4 wurden vom Optimizer zwei optimierte Kanaltypen erstellt. Die Liste der Policys wurde aus Gründen der Übersichtlichekeit ausgelassen. In der aktuellen Version sind die optimierten Kanäle Klassenvariablen, wie in den Zeilen 18 bis 20 sichtbar ist.  Diese Variablen werden in der Initialisierungsliste der Konstruktors (Zeile 9) erstellt und erhalten über den \texttt{PubSubController} Zugriff auf das Netzwerk. Es ist offensichtlich, dass der Optimierer verschiedene Netzwerke mit den einzelnen Instanzen der Klasse \texttt{Channel} verbinden kann.

\lstinputlisting[caption={Ausschnitt der Klasse \texttt{PubSubSystemImpl}}, label=lst:pubsubsystemimpl, float=h!t]{listings/pubsubsystemimpl.h}


\subsubsection*{\texttt{PubSubSystem.cpp}}
Die Klasse \texttt{PubSubSystemImpl} und damit auch die optimierten Kanäle werden erst im Optimierungsschritt erzeugt. Die Klasse \texttt{PubSubSystem} hat jedoch direkten Zugriff auf in \texttt{PubSubSystemImpl} enthaltenen Kanäle und muss daher ebenfalls im Optimerungsschritt angepasst werden. \Fref[plain]{lst:pubsubsystem} zeigt mit der Methode \texttt{getChannelHandle} einen Ausschnitt des generierten Codes. \texttt{impl\_} ist eine Klassenvariable auf eine Instanz von \texttt{PubSubSystemImpl}. Das Codebeispiel zeigt, das \texttt{PubSubSystem} als Dispatcher arbeitet und die entsprechenden Kanäle anhand der \texttt{ChannelList} auswählt. Dieser Zugriff ist nur möglich, da \texttt{PubSubSystem} als \texttt{friend} der Klasse \texttt{PubSubSystemImpl} (\Fref[plain]{lst:pubsubsystemimpl} Zeile 7) definiert ist.

\lstinputlisting[caption={Ausschnitt des zu generierenden Abschnittes von \texttt{PubSubSystem.cpp}}, label=lst:pubsubsystem, float=!t]{listings/pubsubsystem.cpp}

\section{Implementierungsdetails der Klasse \texttt{Channel}}
\label{chap:impl:channel}
In diesem Abschnitt werden Ausschnitte der Klasse \texttt{Channel} gezeigt um damit verschiedene Implementierungsdetails aufzuzeigen. Als erstes wird die Umsetzung eines Teils des Verarbeitungsmodells beim Empfang einer \emph{Publish}-Nachricht in \texttt{deliver} gezeigt. \Fref[plain]{fig:processing_deliver_ovpd} zeigt diesen Teil des Verarbeitungsmodells als Ausschnitt und in \Fref{lst:ovpd} ist der zugehörige Codeausschnitt der Klasse \texttt{Channel} dargestellt. In der Zeile 2 werden die Verwaltungsinformationen der Nachricht ausgelesen und in einem \texttt{Header} gespeichert.  Die Schleife -- beginnend in Zeile 6 -- zeigt die weitere Verarbeitung der \emph{freigegebenen} Nachrichten bishin zur Übergabe an die Applikation (Zeile 16). Die gewählte Strategie der Policy \emph{Reihenfolge} -- hier über \texttt{Order} ansprechbar -- bekommt die zu bearbeitende Nachricht mitsamt einiger Verwaltungsinformationen übergeben. Der Rückgabewert dieser Methode ist eine Liste von \emph{freigegebenen} Nachrichten zusammen mit den weiter benötigten Verwaltunsginformationen. Anhand diesen, kann die Strategie zur \emph{Validierung} -- über \texttt{Validity} ansprechbar -- in Zeile 13 prüfen ob die Nachricht gültig ist oder verworfen werden muss (Zeile 14). Als letzter Schritt vor der Übergabe der Nachricht an die Applikation, kann in Zeile 15, der Event oder dessen Auswirkung auf die virtuelle Welt abgespeichert werden.

\lstinputlisting[caption={Verarbeitung einer \emph{Publish}-Nachricht in \texttt{deliver}}, label=lst:ovpd, float=t]{listings/ovpd.cpp}

\begin{figure}[htbp]
\centering
%\resizebox{\textwidth}{!}{%
\includegraphics{grafics/processing_deliver_ovpd.pdf}%}
\caption{Ausschnitt der Nachrichtenverarbeitung in \texttt{deliver}}
\label{fig:processing_deliver_ovpd}
\end{figure}

Das nächste Beispiel zeigt in \Fref{lst:publish} die Erstellung einer \emph{Publish}-Nachricht und allgemein den Versand einer Nachricht. In der Methode \texttt{publish} werden die Verwaltungsinformationen aller Strategien abgefragt. Der \emph{Filterung} wird zusätzlich die Nachricht übergeben, anhand der die gewählte Strategie Metadaten zu Filterung erzeugen kann. In Zeile 10 wird der erstellte \texttt{Header} und die Nachricht an \texttt{sendMessages} unter Angabe des Types (\texttt{PUBLISH}) übergeben.\\
Im ersten Abschnitt dieser Methode (Zeilen 14 bis 17) wird der generische Nachrichtentyp in eine Nachrichtennummer des Netzwerks umgewandelt. Im zweiten Abschnitt (Zeilen 19 bis 24) wird die eigentliche Nachricht erstellt und verschlüsselt.\\
Für den Fall einer \emph{Publish}-Nachricht wird im dritten Abschnitt (Zeilen 26 bis 33) ermittelt ob die Verteilungsstrategie eine Filterung der Empfängerliste wünscht.\\
Im letzten Abschnitt dieser Methode (Zeilen 35 bis 39), wird die Liste der Empfänger von der Verteilungsstrategie errechnet und in Zeile 38 -- nach erfolgter Prüfung -- über das Netzwerk versandt.

\lstinputlisting[caption={Versenden einer \emph{Publish}-Nachricht}, label=lst:publish, float=t]{listings/publish.cpp}

Das nächste \Fref{lst:header} zeigt Auszüge der, an die gewählten Strategien angepasste, Klasse \texttt{Header}. Diese ist innerhalb der Klasse \texttt{Channel} definiert und hat damit Zugriff auf die, der Klasse \texttt{Channel} übergebenen, gewählten Strategien. Die Größe des Feldes für Verwaltungsinformationen wird anhand der Werte der einzelnen Strategien berechnet. \texttt{KeySize} ist die Länge der Stringrepräsentation eines Schlüssels im Netzwerk und wird der Klasse \texttt{Channel} als Policy übergeben. Jede Strategie \emph{muss} über das Feld \texttt{size} zur Übersetzungszeit angeben, wie viele \texttt{chars} sie zur Speicherung ihrer Metadaten pro Nachricht benötigt. Damit kann die benötigte Größe in Zeile 11 berechnet werden und ein entsprechendes \texttt{char[]} in Zeile 13 erstellt werden. Das restliche Codebespiel zeigt einige der Zugriffsmethoden auf die einzelnen Metadaten. Auch hier ist ersichtlich, dass diese bereits zur Übersetzungszeit mit exakten Werten erstellt werden und somit effektiv gegen Pufferüberlauf schützen. 

\lstinputlisting[caption={An die gewählten Strategien angepasste Klasse \texttt{Header}}, label=lst:header, float=t]{listings/header.h}


%\chapter{Related Work}
\label{chap:related}

IPMulticast ist das Beste, Application Multicast ist toll, aber OverlayMulticast ist besser. OverlayMulticast: explizites Erstellen von RoutingKnoten. Widerspricht aber unserem Ansatz, so wenig "zentrale Server" wie möglich zu haben.\cite{Lao2005Comparative}

\section*{NTree}
\label{chap:related:ntree}
NTree teilt die Spielwelt in Regionen auf. Spieler (Knoten) kommunizieren zuerst nur innerhalb dieser Region und daher kann Kommunikationsaufwand eingespart werden.

\cite{GauthierDickey2005Using}

\paragraph{Arbeitsweise}
Über die Position in der virtuellen Welt werden Knoten in Regionen eingeteilt. Im Paper wird dies exemplarisch durch einen Quadtree aufgezeigt. Pro Region gibt es einen \emph{Leader} der von den Knoten innerhalb der Region eigenständig ermittelt wird. Alle Knoten innerhalb einer Region sind miteinander verbunden. Der Leader ist somit Knoten in einer größeren Region und kann über seine Nachbarn (die selbst Leader ihrer Region sind) mit nahezu allen anderen Knoten im Verbund kommunizieren.

Ein Event E hat einen gewissen Wirkradius.
Löst ein Knoten ein Event aus, so sendet er dieses an all seine Nachbarn. Der Leader einer Region entscheidet, ob das Event einen größeren Wirkradius als die Ausdehnung der Region hat und leitet das Event an seine Nachbarn weiter. Diese geben das Event - falls dies nötig ist - wiederum an ihre eigene Region weiter oder senden es nach oben weiter.

Regionen haben keine statische Größe sondern sind durch die Anzahl der in ihr befindlichen Knoten bestimmt. So werden Regionen gesplittet bzw. zusammengefasst, wenn sich die Anzahl der Knoten ändert beziehungsweise bestimmten Grenzen nähert. Die Ausdehnung einer Region muss durch den Leader berechnet werden.

Da innerhalb einer Region wenig Knoten vorhanden sind, können die Events innerhalb der Region synchronisiert werden. Müssen außerhalb der Region Events synchronisiert werden, so kann das TimeWarp-Verfahren benutzt werden.

\missing{Link auf das TimeWarp-Verfahren!}

\paragraph{Fehlerfall}
Über Heartbeat-Nachrichten wird das Vorhandensein von Knoten abgeprüft. Fällt der Leader einer Region aus, so gibt es zwei Möglichkeiten. Beide Möglichkeiten bedingen, dass einige Knoten einer Region auch den Leader der übergeordneten Region als Nachbar kennen müssen und umgekehrt.

\begin{itemize}
\item Die Knoten der Region ernennen einen neuen Leader aus ihren Reihen. Diese Entscheidung wird an den übergeordneten Leader übermittelt.
\item Der übergeordnete Leader erkennt das Fehlen zuerst und stößt die Leaderwahl an.
\end{itemize}

\paragraph{Unschönheiten}
Knoten in einer Region müssen nicht notwendigerweise benachbart sein und die Kommunikation erfolgt wieder über unbeteiligte Rechner.



\section*{minTCO}
\label{chap:related:mintco}
\cite{Chockler2007Constructing}


\section*{Vivaldi}
\label{chap:related:vivaldi}
Vivaldi \cite{Dabek2004Vivaldi} ist ein dezentrales System, in dem sich Knoten des Netzwerkes Koordinaten so wählen, dass der Abstand zwischen zwei Koordinaten der Latenz zwischen diesen beiden Knoten entspricht.

\paragraph{Arbeitsweise}
Neue Knoten im Netz wählen sich zufällig Koordinaten. Zu allen Nachbarn werden nun Ping-Nachrichten zur Messung der Laufzeit verschickt. Die Koordinaten der einzelnen Knoten werden dabei übertragen. So kann nun die eigene Koordinate anhand des Abstandes angepasst werden.

Das Paper verwendet hier eine Analogie zu gespannten Federn. Alle Knoten sind über Federn miteinander verbunden und suchen nun das gemeinsame Optimum der Federspannung.


\section{VON}
\label{chap:related:von}
\ac{von} ist ein integriertes System aus Overlaynetzwerk mit Publish/Subscribe-System und zielt direkt auf virtuelle Umgebungen wie \ac{mmog} ab \cite{Hu2006VON}. Jeder Knoten hat demnach eine Koordinate im Raum\footnote{z.B. Spielewelt} und eine \ac{aoi}. Die Welt werden mittels Voronoi-Diagrammen in disjunkte Regionen eingeteilt und somit gibt es für jeden Knoten eine Region. Ein Knoten hat direkte und angrenzende Nachbarn. Bewegt sich der Knoten im Raum, so verändert sich die Aufteilung der Regionen.

\ac{vast} \cite{Backhaus2007Voronoibased} greift das Konzept von \ac{von} auf und testet eine Implementierung auf OpenSIM \cite{Baumgart2007OverSim}.


\chapter{Zusammenfassung und Ausblick} 
\label{chap:zus}

Diese Arbeit befasste sich mit der Konzeption und prototypischen Implementierung eines Frameworks zur optimierten Verteilung von Events in Publish/Subscribe-System auf einem strukturiertem p2p-Overlay Netzwerk.\\
Nach der Einführung von \ac{m2etis} wurden die Grundlagen zu Overlay- und p2p-Netzwerken gelegt. Eine Kapitel über betrügerisches Verhalten und dessen Verhinderung runden die Einführung ab.\\
Im zweiten Abschnitt der Arbeit wurden die drei Overlay-Netzwerke Chord, Pastry/Tapestry und CAN vorgestellt und miteinander verglichen um einen auf die identifizierte Anforderungen passenden Netzwerktyp für \ac{m2etis} zu finden. Mit Chimera hat sich ein freies Netzwerksystem gefunden, dass das gut geeignete Konzept von Pastry/Tapestry umsetzt.\\
Auf die Evaluation folgte die Konzeption des Frameworks. Die in \cite{Fischer2010Event} identifizierten orthogonalen Dimensionen sind durch sieben Policies abgedeckt, welche Auswirkungen auf die Nachrichtenverteilung und damit die Bearbeitung der Events haben. Im Abschnitt zum Verarbeitungsmodell wird die Reihenfolge der verschiedenen Policies und ihr Zusammenspiel bei der Verarbeitung der Nachrichten erklärt. Die beispielhafte Anwendung verschiedenster Verteilungsstrategien zeigen die Mächtigkeit und Offenheit des Modells viele Implmentierungen zu unterstützen und garantiert damit eine Optmierung der Eventverteilung.\\
Eine Einführung in Template Meta-Programming stellt einige die zur Reduzierung von Overhead zur Laufzeit genutzten Techniken vor. Einblicke in die Implementierung des Frameworks zeigen die Abstraktionsebenen des Frameworks auf und komplettieren dieses Kapitel.


\subsection*{Ausblick}
\begin{itemize}
\item verschiedene Strategien
\item Kostenmodell
\item semantisches Modell
\item DSL zur Beschreibung der Eventtypen
\item Netzwerksimulator OverSim \cite{Baumgart2007OverSim}
\item Tuning Kotenmodell
\item Tuning Parameter
\item Test Optimizer
\end{itemize}

Als weitere Arbeiten in diesem Gebiet steht die Entwicklung des Optimierers und der Kostenmodelle an nächster Stelle. Eine domänenspezifische Sprache zur semantischen Beschreibung der Typen rundet diesen ab. Zur Parametrisierung der verschiedenen Strategien kann das System vorübergehend an ein simuliertes Netzwerk wie OverSim \cite{Baumgart2007OverSim} angeschlossen werden. Weitere Algorithmen müssen in das System eingebracht und ausgemessen werden. Durch die flexible Architektur ist es möglich mehrere logische Netzwerke zu nutzen und bestimmten Eventtypen zuzuordnen.



\appendix

\section{Moderne Entwicklung mit C++}
\label{chap:impl_tmp}
\acf{tmp}\index{Template Meta-Programming} und policy-based Design\index{policy-based Design} \cite{Alexandrescu2001Modern} sind Paradigmen der Softwareentwicklung in C++ und bedienen sich der dort verfügbaren \emph{Templates}. Policy-based Design baut in der hier erklärten Variante auf dem turingvollständigem \ac{tmp} auf.

\subsection{Templates}
Templates dienen der Generalisierung von Code und werden zur Übersetzungszeit durch den Compiler ausgewertet. Templates sind ein integraler Bestandteil der \ac{stl}\footnote{Siehe auch \url{http://www.cplusplus.com/reference/stl/}} und dienen dort vor Allem zum Erstellen von abstrakten Containerklassen wie zum Beispiel \texttt{std::vector}, \texttt{std::map} oder \texttt{std::set}. Basierend auf dem Konzept des Zugriffs über \texttt{Iteratoren}, sind viele aus der Funktionalen Programmierung bekannte Funktionen implementiert, die ebenfalls durch Templates generisch gehalten sind.

Einzelne Funktionen oder komplette Klassen könnten als Template realisiert werden. Zusätzlich der generischen Implementierung können auch spezielle Implementierungen, sog. Spezialisierungen, angegeben werden können. Funktionstemplates lassen sich nur total, Klassentemplates auch partiell spezialisieren. Eine partielle Spezialisierung kann bei mehreren Templateparametern erfolgen. Ungenutzte Templates werden vom Compiler nicht instantiiert, dies bedeutet dass kein ungenutzter Code im fertigem Kompilat vorhanden ist. Dies wiederum verringert die Größe von ausführbaren Dateien oder dynamischen Bibliotheken. Jedoch muss erwähnt werden, dass jede Instanz eines Templates eigenen Code generiert. Dieser lässt sich - dank des Wissens zur Übersetzungszeit - besser optimieren als zur Laufzeit.

In diesem und den folgenden Beispielen soll ein Einblick in die neue Methodik gezeigt werden, daher wird nicht auf eine optimierte Parameterübergabe oder Ähnliches geachtet. \Fref{lst:tmp_easy_spez} zeigt ein einfaches Template. Der Funktion \texttt{min} wird als Templateparameter der Typ der Parameter übergeben. Es wird erwartet, dass die übergebenen Typen den Operator $<$ überladen haben. Für den Typ \texttt{Pair} ist das Template total spezialisiert und eine separate Implementierung wurde angegeben. Zeilen 18 und 19 zeigen die einzelnen Funktionsaufrufe.

\lstinputlisting[caption={einfaches Funktionstemplate}, label=lst:tmp_easy_spez, float=t]{listings/tmp_easy_spez.cpp}

Ebenso können Templates geschachtelt übergeben werden, wie es \Fref{lst:tmp_tmp} zeigt. Hier wird dem \texttt{WidgetFactory} ein Template übergeben mit dessen Hilfe er Objekte des Typs \texttt{Widget} erstellen kann. Beim Aufruf in Zeile 19 muss der Entwickler dennoch Typ angeben obwohl eine \texttt{WidgetFactory} nur Objekte des Typs \texttt{Widget} erzeugt, wie es der Name impliziert.

\lstinputlisting[caption={einfaches Klassentemplate}, label=lst:tmp_tmp, float=t]{listings/tmp_tmp.cpp}

\Fref{lst:tmp_tmp_2} zeigt die Lösung mit \enquote{Template Template} Parametern. Hier wird der \texttt{WidgetFactory} ein Template übergeben. Dieses Template ist noch nicht spezifiziert und benötigt weiterhin den Typ des zu erstellenden Objektes. Der Unterschied zum vorherigem Listing ist in Zeile 5 und am Aufruf in Zeile 10 zu sehen. \texttt{WidgetFactory} übergibt den Typ an das Template. Der Entwickler ist nicht mehr für die Auswahl des korrekten Typs zuständig.

\lstinputlisting[caption={verbessertes Klassentemplate mit Template Template}, label=lst:tmp_tmp_2, float=!t]{listings/tmp_tmp_2.cpp}

Mit Templates lassen sich auch Entscheidungen anhand der übergebenen Templateparamter treffen. In \Fref{lst:tmp_cond} wird anhand des Parameters \texttt{isSingleton} entschieden, ob das neue Objekt mit \texttt{new} oder per Aufruf der Methode \texttt{T::getInstance()} erzeugt werden kann.

\lstinputlisting[caption={Entscheidung via Templates}, label=lst:tmp_cond, float=t]{listings/tmp_cond.cpp}

Der Compiler prüft dabei alle Templates syntaktisch. Obwohl \texttt{isSingleton} mit \texttt{false} instantiert wird, kann der Code nicht kompiliert werden, denn die Methode \texttt{getInstance} wird von \emph{Foo} nicht angeboten. Um dieses Problem zu umgehen, beschreibt Alexandrescu in \cite{Alexandrescu2001Modern} das Hilfskonstrukt \texttt{Int2Type} das im nächsten Kapitel eingeführt wird.

\subsection{Template Meta-Programming}
Template Meta-Programming arbeitet mit der vorgestellten Spezialisierung von Templates und wird ebenfalls vom Compiler zum Übersetzungszeitpunkt ausgewählt und basiert rein auf Typinformation.

\ac{tmp} kann dazu genutzt werden, um die Problematik im vorangegangenen Kapitel zu lösen. \Fref{lst:tmp_cond_2} zeigt die Definition von \texttt{Int2Type}. \texttt{Int2Type} speichert den übergebenen Integerwert und erzeugt daraus pro Wert einen eigenen Typ. Mit überladenen Funktionen kann nun eine Auswahl zur Übersetzungszeit erfolgen. Die statische Methode \texttt{create} (Zeile 8) ruft nun abhängig vom Templateparameter \texttt{isSingleton} die überladene Methode \texttt{create\_impl} auf. Im ersten Fall (Zeile 14) nimmt diese ein \texttt{Int2Type$<$true$>$}, im zweiten Falle (Zeile 19) nur \texttt{Int2Type$<$false$>$}. Da der Compiler nur instantiierte Templates prüfen muss, gibt es nun keine Übersetzungsfehler mehr. Der Entwickler muss seinen Aufruf nicht anpassen.

\lstinputlisting[caption={Entscheidung via Templates zur Übersetzungzeit}, label=lst:tmp_cond_2, float=t]{listings/tmp_cond_2.cpp}

Durch diese Technik können Entscheidungen zur Übersetzungszeit getroffen werden. So können zum Beispiel auch komplexe Berechnungen zur Übersetzungszeit ein- beziehungsweise ausgeschaltet werden. Der zusätzliche Code (zum Beispiel Methodenaufruf) kann durch den Compiler leicht optimiert werden.

\subsection{Policy-based Design}
Policy-based Design kombiniert Vererbung und Templates und ermöglicht es generischen Code zu entwickeln. Hierbei lassen sich sehr gut orthogonale Aspekte ausnutzen, wie Alexandrescu am Beispiel eines SmartPointers erläutert.

In \Fref{lst:tmp_pbd} ist ein einfaches Beispiel mit einer Ausgabeklasse (\emph{Ausgabe}) beschrieben, die mittels einem \emph{Decorator} und einem \emph{Printer} parametrisiert wird. Ausgabe leitet von diesen beiden Typen ab und nutzt deren Funktionalität. Dadurch erzeugt der Compiler implizit Anforderungen die diese Parameter zu erfüllen haben. Jeder Decorator muss eine Methode \texttt{decorate} anbieten, welche einen \emph{const string\&} als Parameter nimmt. Jeder Printer muss eine Methode \texttt{print} anbieten, welche den Rückgabetyp von Decorator::decorate als Parameter nimmt. Diese Kontrakte werden vom Compiler zur Übersetzungszeit in Zeile 24 geprüft. Weitere Varianten für Decorator und Printer sind trivial.

Ausgabe wird in diesem Zusammenhang als \enquote{Host-Klasse}, Dec\_A und Printer als \enquote{Policy-Klasse} bezeichnet. Da die Host-Klasse von ihren Policy-Klassen ableitet, kann die Host-Klasse in eine Policy-Klasse gecastet werden. So können nützliche Informationen abgefragt oder speziell angepasste Methoden der Policy-Klasse aufgerufen werden. Der Nutzer einer solchen mit Policies angereicherten Klasse bekommt dadurch sehr viele Freiheiten. Mögliche Sonderfälle spricht Aleandrescu an und bietet angepasste Lösungen.

\lstinputlisting[caption={Beispiel für policy-based Design}, label=lst:tmp_pbd, float=t]{listings/tmp_pbd.cpp}

Diese Möglichkeit der Parametrisierung vom Klassen gibt es prinzipell auch zur Laufzeit. Hier wird jede Policy durch einen Pointer auf eine abstrakte Basisklasse abgebildet. Jedoch hat dies einige Nachteile zur Laufzeit. Zum einen sind die Interfaces der Policy-Klassen festgeschrieben und die Host-Klasse leitet nicht von diesen ab. Zusätzlich sind virtuelle Funktionsaufrufe \enquote{teuer}, da diese in der sogenannten \enquote{VTable} zur Laufzeit abgeprüft werden müssen. Diese Überprüfung kann durch den Compiler nicht entfernt oder optimiert werden. 


In dieser Arbeit kommt \ac{tmp} im Zusammenspiel mit policy-based Design bei der Entwicklung des einzelnen Kanals zur Anwendung. Durch die Optimierungen zur Übersetzungszeit wird schlanker Code ohne unnötigen Ballast zur Laufzeit erzeugt.


\chapter{Distributed Hashtable (DHT)}
\label{chap:dht}

Das System der \acf{dht} ist nach \cite{Wehrle2005} eine Lösung des \emph{lookup}-Problems\footnote{\enquote{Wo speichere beziehungsweise finde ich Datensatz X?}} in verteilten Systemen und besticht dabei durch seine Einfachheit, Robustheit und Skalierbarkeit. \acp{dht} können einerseits zur Speicherung von Datensätzen nach dem \emph{(key,value)}-Ansatz genutzt werden oder zum Routen von Nachrichten durch das Netzwerk. Beide Varianten schließen sich nicht gegenseitig aus. Für das System ist eine Hashfunktion $h$ definiert, so dass alle Eingaben eindeutig auf einen Wert aus dem \emph{Schlüsselraum} $S$ abgebildet werden: $h(x) \rightarrow y; y \in S$. Eine \emph{segmentierte} Hashfunktion bildet den Eingabewert auf $n$ eindeutige Schlüssel ab. Dabei ist der Schlüsselraum in $n$ disjunkte Bereiche aufgeteilt: $h_2(x) \rightarrow y,w; y \in S_1, w \in S_2, S := S_1 \cup S_2$. Die Wertemenge des Schlüsselraumes besteht dabei meist aus einem Intervall aus Ganzzahlen, beispielsweise von $0$ bis $2^{160}-1$. Die Hashfunktion bildet dabei die Eingaben möglichst auf den ganzen Wertebereich ab.

Jeder Knoten im Netzwerk bekommt über diese Hashfunktion einen eindeutigen Schlüssel zugeordnet. Eine mögliche Eingabe der Hashfunktion ist zum Bespiel eine Kombination aus Hostnamen und Port. Einem Knoten wird -- je nach Netzwerk -- ein gewisser Bereich aus dem Schlüsselraum zugeordnet. Weiterhin hat jeder Knoten eine gewisse Kenntnis von andern Knoten im Netzwerk -- seine Nachbarschaft. Diese erlangt er durch netzwerkabhängige Verfahren beim Eintritt in das System\footnote{Beispiele sind in \Fref{chap:evaluation_p2p} bei der Vorstellung der verschiedenen Netzwerke Chord, Pastry und CAN gegeben.}. CAN benötigt eine an die Dimensionen angepasste segmentierte Hashfunktion zur Berechnung der Koordinaten.

Wird das System zur Übertragung von Nachrichten genutzt, so besteht dessen minimale Schnittstelle aus den Methoden \texttt{route(key, message)} zum Senden und \texttt{receive(message)} zum Empfangen von Nachrichten\footnote{Siehe \Fref{chap:grundlagen:api}}. Das System kann die Nachricht in jedem Schritt zu einem Knoten senden, der näher -- meist numerisch gesehen -- am Empfänger ist.\\
Wird das System als verteilter Datenspeicher genutzt, so wird für jeden Datensatz dieselbe Hashfunktion genutzt. Beispielweise wird der Dateiname oder der Dateiinhalt an die Hashfunktion übergeben. Damit kann für jeden Datensatz ein eindeutiger Schlüssel berechnet und damit der zuständige Knoten bestimmt werden. Dieser speichert den Datensatz oder hält Informationen bereit, um auf den Datensatz zugreifen zu können. Die minimale Schnittstelle eines solchen Systems beschränkt sich auf die Methoden \texttt{put(key, value)} zum Speichern eines Datensatzes und \texttt{get(key)} zur Abfrage.


%\chapter{Voronoi-Diagramme und Delaunay-Triangulation}
\label{chap:voronoi}


\begin{figure}[htbp]
\centering
\includegraphics{grafics/voronoi_grund_aus_aurenhammer.pdf}
\caption{Voronoi-Diagramm für acht Punkte (aus \cite{Aurenhammer1991Voronoi})}
\label{fig:voronoi_grund}
\end{figure}





\begin{figure}[htbp]
\centering
\subfigure[Dualität von Voronoi-Diagramm und De\-launay-Triangulation (aus \cite{Aurenhammer1991Voronoi})]{
	\includegraphics{grafics/voronoi_dual_aus_aurenhammer.pdf}
	\label{fig:voronoi_dual}
}
\subfigure[Bezug eines Voronoi-Diagrammes zur konvexen Hülle (aus \cite{Aurenhammer1991Voronoi})]{
	\includegraphics{grafics/voronoi_konvex_aus_aurenhammer.pdf}
	\label{fig:voronoi_konvex}
}
\end{figure}


\cite{Aurenhammer1991Voronoi}

\cite{Fortune1987, Dwyer1987} %Erzeugen von Voronoi-Diagrammen (Sweepline und Co)



% CONTENT END

\cleardoublepage

\automark[chapter]{section}

% Index soll Stichwortverzeichnis heissen
\newpage
\renewcommand{\indexname}{Stichwortverzeichnis}
	 
% Stichwortverzeichnis soll im Inhaltsverzeichnis auftauchen
\phantomsection
\addcontentsline{toc}{chapter}{\indexname}

\markleft{\indexname}
\markright{\indexname}
\printindex

\automark[chapter]{section}

\bibliographystyle{alphadin}
\bibliography{bib/literatur}
%%% Wenn folgende Zeile nicht auskommentiert wird, steht alles im Litverzeichnis
%\nocite{*}
\end{document}
